%% The report to gather all sections into, set formatting and compile.

\documentclass[a4paper, 11pt, hidelinks]{article} %hidelinks removes the red boxes around hyperlinks
%TODO prevent word spilling over lines
%%%%%%%%%% Packages %%%%%%%%%%
\usepackage{amsmath}
\usepackage{epstopdf} % to allow use of crest in title page
\usepackage{graphicx}
\usepackage{float}
\usepackage[mathlines]{lineno}%line numbers
\usepackage{textcomp} % for degree symbol
%biblatex
%			compiler		style				max intext names	maxinbib names  
\usepackage[backend=biber, style=authoryear-comp, maxcitenames=2, maxbibnames=10,
%sirnames same add initial  names same add next author
uniquename=false, 			uniquelist=false, 
%no urls	no isbn	
url=false, isbn=false]{biblatex}
\DeclareDelimFormat{nameyeardelim}{\addcomma\space} % get commas between author and year

\usepackage{hyperref} % for URLs
\usepackage[table, dvipsnames]{xcolor}
%%%%%%%%%% Custome Commands %%%%%%%%%%
\newcommand{\crest}{\includegraphics[width = 4cm, keepaspectratio]{../images/IC_Crest.eps}} % Imperial crest
\newcommand\wordcount{%% The report to gather all sections into, set formatting and compile.

\documentclass[a4paper, 11pt, hidelinks]{article} %hidelinks removes the red boxes around hyperlinks

%%%%%%%%%% Packages %%%%%%%%%%
\usepackage{amsmath}
\usepackage{epstopdf} % to allow use of crest in title page
\usepackage{graphicx}
\usepackage[mathlines]{lineno}%line numbers
\usepackage{textcomp} % for degree symbol
%biblatex
%			compiler		style				max intext names	maxinbib names  
\usepackage[backend=biber, style=authoryear-comp, maxcitenames=2, maxbibnames=10,
%sirnames same add initial  names same add next author
uniquename=false, 			uniquelist=false, 
%no urls	no isbn	
url=false, isbn=false]{biblatex}
\DeclareDelimFormat{nameyeardelim}{\addcomma\space} % get commas between author and year

\usepackage{hyperref} % for URLs
\usepackage[table, dvipsnames]{xcolor}
%%%%%%%%%% Custome Commands %%%%%%%%%%
\newcommand{\crest}{\includegraphics[width = 4cm, keepaspectratio]{../images/IC_Crest.eps}} % Imperial crest
\newcommand\wordcount{%% The report to gather all sections into, set formatting and compile.

\documentclass[a4paper, 11pt, hidelinks]{article} %hidelinks removes the red boxes around hyperlinks

%%%%%%%%%% Packages %%%%%%%%%%
\usepackage{amsmath}
\usepackage{epstopdf} % to allow use of crest in title page
\usepackage{graphicx}
\usepackage[mathlines]{lineno}%line numbers
\usepackage{textcomp} % for degree symbol
%biblatex
%			compiler		style				max intext names	maxinbib names  
\usepackage[backend=biber, style=authoryear-comp, maxcitenames=2, maxbibnames=10,
%sirnames same add initial  names same add next author
uniquename=false, 			uniquelist=false, 
%no urls	no isbn	
url=false, isbn=false]{biblatex}
\DeclareDelimFormat{nameyeardelim}{\addcomma\space} % get commas between author and year

\usepackage{hyperref} % for URLs
\usepackage[table, dvipsnames]{xcolor}
%%%%%%%%%% Custome Commands %%%%%%%%%%
\newcommand{\crest}{\includegraphics[width = 4cm, keepaspectratio]{../images/IC_Crest.eps}} % Imperial crest
\newcommand\wordcount{%% The report to gather all sections into, set formatting and compile.

\documentclass[a4paper, 11pt, hidelinks]{article} %hidelinks removes the red boxes around hyperlinks

%%%%%%%%%% Packages %%%%%%%%%%
\usepackage{amsmath}
\usepackage{epstopdf} % to allow use of crest in title page
\usepackage{graphicx}
\usepackage[mathlines]{lineno}%line numbers
\usepackage{textcomp} % for degree symbol
%biblatex
%			compiler		style				max intext names	maxinbib names  
\usepackage[backend=biber, style=authoryear-comp, maxcitenames=2, maxbibnames=10,
%sirnames same add initial  names same add next author
uniquename=false, 			uniquelist=false, 
%no urls	no isbn	
url=false, isbn=false]{biblatex}
\DeclareDelimFormat{nameyeardelim}{\addcomma\space} % get commas between author and year

\usepackage{hyperref} % for URLs
\usepackage[table, dvipsnames]{xcolor}
%%%%%%%%%% Custome Commands %%%%%%%%%%
\newcommand{\crest}{\includegraphics[width = 4cm, keepaspectratio]{../images/IC_Crest.eps}} % Imperial crest
\newcommand\wordcount{\input{Report.sum}} % to include word count 

%%%%%%%%%% Formatting %%%%%%%%%%
\usepackage[margin=2cm]{geometry} % margins of 2cm
\linespread{1.5} %1.5 spacing
%\renewcommand{\familydefault}{\sfdefault} % set font to arial clone (helvet)

\usepackage[compact]{titlesec} % reduce spacing bewteen section titles

%%%%%%%%%% Bibliography %%%%%%%%%%
\addbibresource{../Masters_Thesis.bib}

%commands 

%%%%%%%%%% Document %%%%%%%%%%
\begin{document}
	%%%%%%%%%% Title Page %%%%%%%%%%
	\include{title_page}
	

	%%%%%%%%%% Declaration %%%%%%%%%%
	\section*{Declaration}
	I declare this project as my own work.  The model presented here was developed in conjunction with my supervisor, Dr. Samraat Pawar, and Ph.D. students Tom Clegg and Olivia Morris.  I was responsible for any simulations and data presentation.\newline
	%% word count
	\textbf{Word Count: \wordcount}

	\newpage
	
	%%%%%%%%%% Abstract %%%%%%%%%%
	\section*{Abstract}
	\linenumbers
%	Size is essential to reproductive output. By extension understanding growth, determines size allows understanding of reproductive output. 
	 
	With recent results showing that reproduction in fish scales hyperallometrically there is a need to update growth OGMs to reflect this fact.  Current OGMs assume optimal intake, an assumption which is not always reflected in the field.  This study develops an energy intake focused approach to explaining growth, an area which has not been covered within current literature, and shows that hyperallometric scaling of reproductive output arises when allowing for variable reproductive scaling and maximising for fitness.  The model is applicable to not only fish, but any animals taxon with some simple parameter adjustments.  I offer direction for improvements and areas to be developed in order to allow the model to be applicable to any temperature range.
	\vspace*{0.5 cm}
	\newline
%	TODO pieces to include:
%	
%	show the effects of the exponent and why it is important that we use the correct one moving forward.
%	give direction as to what needs to be done going forward.  what areas need work
	%TODO KEywords
	\textbf{Keywords:}\\
	allometry; functional response; metabolic theory; growth; intake; life history; metabolism; reproduction; reproductive output; supply

	
	
	\nolinenumbers
	%%%%%%%%%% Acknowledgements %%%%%%%%%%
	\include{Acknowledgements}
	
	%%%%%%%%%% Table of Contents %%%%%%%%%%
	\tableofcontents
	\newpage

	%%%%%%%%%% Introduction %%%%%%%%%%
\section{Introduction}
	\linenumbers
	% putting OG first
	Growth is a gateway to understanding how many biological rates change throughout ontogeny.  This is because so many biological rates scale with mass \parencite{Kleiber1932}.  To date many models have been developed to predict and describe the growth of an organism throughout its lifetime.  The three main approaches used are the von Bertalanffy model, the dynamic energy budget (DEB) model and the ontogenetic growth model (OGM), which will be the focus of this study \parencite{Putter1918, vonBertalanffy1938, Kooijman1986, West2001}.  All of these are energetic based models with varying assumptions, key among which is the scaling of resource supply and metabolic rate with mass.   In OGMs supply is thought of as being optimal at all times which leads to the assumption that intake scales to the power of 0.75.  Indeed while under optimal conditions this may be true, however, it neglects that this situation is thought to rarely occur in the field \parencite{Pawar2012}.
%	Additionally, growth model approaches have made various assumptions about the scaling of metabolic rates.  OGMs assume that scaling is isometric with mass, i.e. scaling exponent is 1.  This is contrary to what has been thought 
	This study takes the novel approach of using a mass specific functional response and assimilation efficiency to describe how intake changes both throughout ontogeny and varying levels of resource availability. While the methodology used here can be used for any animal taxon, in order to investigate growth, a taxon much be chosen to parametrise metabolic rate around.  This study will focus on fish.
		
% TODO bring in OGM growth curve and use to illustate the basic OGM as per samraats comments
	% Why is growth important
	Understanding how organisms grow and what factors play a role in determining growth is not only insightful from the perspective of understanding the world around us, but can also be used to better manage the many fisheries and marine protected areas around the world.  An objective which is becoming increasingly important as the oceans' fish stocks continue to be depleted by over fishing. 
	%VV rephrase this part to hint more towards size than it currently does VV
	This is compounded by global warming which threatens to alter the structure of the ecosystems even if they are not fished and left in their ``natural" state. 
	
	Larger fish produce more offspring than smaller ones.  % this is known in general isometrically
	In fact, it has been shown that larger fish produce far more offspring than the equivalent mass composed of smaller fish,  that is a 2kg fish will produce more offspring than two 1kg fish \parencite{Barneche2018}.
	Larger fish also use energy more efficiently than multiple smaller ones per unit mass.  This is due to them having a lower mass specific metabolic rate \parencite{Peters1983, Kleiber1932, Brown2004}.  
	%As a result a higher proportion of larger fish may actually reduce stress upon the ecosystem from an energetic perspective which could be useful when trying to manage fish stocks. %need citations for this also questionably neccessary as a point and has little backup? may even not be the case at all due to other pressures introduced by larger individuals in a system
	% larger mothers produce larger offspring which may better survive \cite{Barneche2018}
	Additionally, larger mothers produce larger offspring, which are then more likely to survive to adulthood and reproduce \parencite{Hixon2014, Marshall2006}. % double check this statement from barneche and get more backing
	%
	
	It is already known that metabolic rate is dependant on temperature which in turn will lower fish sizes \parencite{Gillooly2001, Brown2004} and with increasing global temperatures, understanding in greater detail how increased metabolic rates may affect growth is useful in population management.
	%
	
%	Growth is in essence a balancing act between how much energy an organism can acquire and the amount of energy required for maintenance, 
%	that is movement, digestion et cetera.  
	
	%power laws -- needs citation or could be counted as common knowledge? \cite{Huxley1932}
	Key to understanding biological rates and their relationship with mass is the concept of power laws. Many biological traits can be described as scaling to the power of some other biological trait.  That is some rate, $Y$, can be expressed for any mass by $Y = Y_0 m^\beta$, where $m$ is mass \parencite{Kleiber1932}.  Power laws can be broadly categorised based on the value of their exponent, $\beta$.  
	Where the exponent equals one the relationship is described as isometric or linear, that is the two traits increase at the same rate.  
	Where the exponent does not equal one, the relationship is said to be allometric.  That is the trait does not increase at the same rate as the trait being compared against.  In cases where the exponent is greater than one, the trait is said to scale super-linearly or hyperallometrically and to scale sub-linearly or hypoallometrically for cases with an exponent of less than one.
	For the purposes of describing growth these relationships are indispensable.   
	
	
%	Traditionally ontogenetic growth models have relied on knowing how large an organism is expected to grow.
%	%	von Bertalanffy / puetter model
%	The von Bertalanffy growth equation relies on knowing the longest a fish can be and the length of the fish at the beginning of measurements \parencite{vonBertalanffy1938, Putter1918}.  From here, for a known growth rate, the length of the fish after a set amount of time has passed can be predicted.
%	
%	%OGMs
%	Another approach has been to try explain growth from an energetic perspective.  That is an individuals growth is determined by the amount of energy available to it.  Once all costs, such as maintenance, have been paid all remaining energy can be invested towards growth.  There are two frameworks which have been built around this concept; dynamic energy budget (DEB) theory % TODO get ref for DEB original/important sources %add note to see discussion for more on DEB?
%	and ontogenetic growth models (OGMs), which will be focused on here.

	%	introduce \cite{West2001}	
	One of the best known examples of an OGM is the model developed by \cite{West2001}.  This model is parametrised around the average energy content of animal tissue and asymptotic mass.  Asymptotic mass being the mass at which growth has essentially stopped due to metabolic cost and energy intake equalling each other (Fig. \ref{scaling_plot}a). The model hinges on the scaling with mass between energy intake (m$^{0.75}$, i.e. allometric sub-linear scaling) and maintenance cost (m$^1$, i.e. isometric linear scaling).  In other words, as mass increases maintenance costs will slowly overtake the intake rate and halt growth.  	
	
	% talk about determinant and indeterminant growth/
	% move to including reproduction with \cite{Charnov2001} mashed with \cite{Hou2008} imporvements briefly
	The framework used by \citeauthor{West2001} (\citeyear{West2001}) was latter developed by \citeauthor{Charnov2001} (\citeyear{Charnov2001}) to take the cost of reproduction into account and allowing the estimation of lifetime production of offspring.  \citeauthor{Hou2008} (\citeyear{Hou2008}) latter developed \citeauthor{West2001}'s model further by expanding maintenance cost to include the cost of feeding and digestion (specific dynamic action), synthesis of new tissue and activity.
	
%	% move to including reproduction with \cite{Charnov2001}
%	The framework used by \citeauthor{West2001} (\citeyear{West2001}) was latter developed by \citeauthor{Charnov2001} (\citeyear{Charnov2001}) to take the cost of reproduction into account and allowing the estimation of lifetime production of offspring.  This was done by adding an additional cost which dedicated some portion of mass to reproductive output, essentially gonadosomatic index.	%like 20% max? in fish very high in  invertabrates \cite{Parker2018}
%	% improved by hou et al \cite{Hou2008} talk about including SDA and energy improvements
%	\citeauthor{West2001}'s (\citeyear{West2001}) model was later expanded upon by \citeauthor{Hou2008} \citeyear{Hou2008} to elaborate on the maintenance cost.  \citeauthor{Hou2008} (\citeyear{Hou2008}) takes both the difference between the cost of maintaining tissue and creating it, a factor that is neglected in \citeauthor{West2001} (\citeyear{West2001}), and the cost of digestion into account.
	
	% begin caveats
	% tautology still present 

	% Discuss allometry and isometry here to highlight what scaling super or sub linearly means
	In the above OGMs intake is assumed to scale sub-linearly to the power of 0.75.  This is due to the assumption that the individuals are consuming at an optimal rate at all times and therefore the only limitation is their ability to make use of that energy.  In this case, intake should theoretically scale to the power of 0.75 (see \cite{West1997}).  However, this is not always the case in the field.  It has been shown that, for non-optimal consumption, steeper scaling can occur \parencite{Pawar2012, Peters1983}. Additionally, OGMs, like many growth and metabolic models, typically use basal or resting metabolic rate to calculate metabolic cost.  This is the minimal metabolic rate of an organism and is typically thought of as the rate of the organism when at relaxed and at rest.  However, it has been shown, once factors such as movement are taken into account, that the scaling becomes steeper \parencite{Weibel2004}.
	% still no accounting for imperfect intake though is somewhat covered by \cite{Hou2011}
	% bring up the caveats/assumptions of the model i.e. isometric metabolism and repro, 0.75 scaling of intake. optimal intake at all times etc. 
	% tautology of OGMs and \cite{Hou2011} close but still issues
	The issue of non-optimal feeding is addressed somewhat by \citeauthor{Hou2011} (\citeyear{Hou2011}).  However, this growth was only investigated as, essentially, a proportion of optimal consumption and does not address a potential change in scaling of intake rate.
	Another limitation of the models used in previous OGMs is the dependence on asymptotic mass.  The models are entirely dependant on the value of optimal intake and asymptotic mass.  All other values, such as metabolic cost, are then derived in relation to these.  However, organisms are not born with some inherent restriction to the size they can attain, at least not energetically.  If there is surplus energy for a given mass the organism should be able to grow.  Relying on asymptotic mass to define the upper bound of attainable mass does not allow for investigation of the mechanisms that underpin asymptotic size in reality. 
	
	% talk about how we assume that evolution settles upon values which optimise for reproduction -- can do this when talking about charnov?
	
	% Justify and appeal my methods	
	With two of the key assumptions of current OGMs, that reproduction and metabolism scale isometrically, not holding in the field \parencite{Barneche2018, Pawar2012, Peters1983} there is a need to take an unexplored approach to modelling fish growth, in particular choosing to focus on developing how intake is described so as to better reflect the real world.  To do this an obvious starting point is to model intake as a functional response \parencite{Holling1959} so as to better reflect real world intake rates in terms of consumed biomass over time.  Non-optimal supply is a currently unexplored area within growth modelling.  This is likely due to the difficulty of directly measuring intake, especially in the field, Perhaps as a result comparatively less is known about consumption.  This leads to the need for proxy values to estimate it, for example nutrient flux \parencite{Schiettekatte2020}, or drawing broad relationships to approximate consumption, as this study will do.
	Changing the manner in which intake is defined also requires changing metabolic cost, because the two are dependent upon each other in current OGMs.  This can be done by defining metabolic rate as a value dependent on current mass rather than asymptotic mass as has been done in OGMs up until this point.  This thought process is more mechanistic since an organism has no concept of ``How large should I grow?", but rather will acquire as much resources as it capable of at its current life stage and size.  Taking this more bottom-up mechanistic approach also allows exploration of what controls how large fish grow to, since as previously mentioned, from an energetic standpoint, an organism can grow indefinitely as long as there is surplus energy available after costs have been paid.  Of course, there are also mechanical and genetic limitations upon organism size, but once size is constrained to what is known to exist, this should not be an issue.  
	Under the assumption that fish have evolved to maximise reproductive output, using the above framework, simulations can be carried out to demonstrate what conditions need to be met in order to achieve hyperallometric scaling of reproduction from an energetic perspective.  This study will answer 1) Under what metabolic scaling exponent can hyperallometric reproduction arise, 2) 
	% describe basic OGM and how i change the intake term to reflect supply
	
%	\cite{Sibly2020} % for metabolic scaling
%	\cite{Marshall2019} non-isometric scaling of repro#

% questions
% what effect does exponent actually have in relation to growth
% what coditions need to be fuflilled to possibly acheive hyperallometry
% 
	\nolinenumbers
	
	%%%%%%%%%% Methods %%%%%%%%%%
\section{Methods}
	\linenumbers
	
	\subsection{Altering OGMs to account for supply}
	In order to address the issue of supply in the context of an OGM, which can be generically described as $dm/dt = gain - loss$, some changes need to be made to the model's terms.  The first is to remove the assumption of asymptotic mass and the reliance of cost upon it.  Within a tradional OGM the gain term ($ a $) and asymptotic mass are used to define the metabolic cost ($ b $).  However, since the assumption of perfect intake is going to be broken, because of the variable supply, this relationship no longer holds.  As such, both intake and metabolic cost need to be redefined.  Additionally, in light of recent work showing that reproduction scales allometrically and not isometrically, the reproductive cost must also be modified from the form used by \citeauthor{Charnov2001} (\citeyear{Charnov2001}) \parencite{Marshall2019, Barneche2018} %TODO check marshal ref data origin
	
	
	\subsubsection{Gain}
	To define intake a natural starting place is the functional response \parencite{Holling1959}.  Functional responses  are used to define how much an organism consumes for a given resource density and is described by the following equation:	
	\begin{equation}
		\label{functional_repsonse}
		f(\cdot) = \frac{a X_r}{1 + a h X_r}
	\end{equation}
	where, $ f(\cdot) $ is the functional response, $ a $ is the search rate, $ h $ is handling time and $ X_r $ is resource density.  
	For a fixed mass and increasing resource density Eq \ref{functional_repsonse} produces a sigmoidal shape with intake eventually reaching an asymptote after some saturating amount of resources is reached.  At lower resource densities, the intake rate is primarily defined by the search rate with higher search rates yielding higher intake rates.  Conversely, at high resource densities, intake rate is approximately equal to the inverse of the handling time ($ h^{-1} $), where lower handling times yield higher intake rates.  
	
	An organism's functional response will not remain constant throughout its life history.  Search rate and handling time are affected by both the organism's mass and how it interacts with its environment \parencite{Pawar2012}.  
%	Within this model mass will be known for all time points since that is one of the quantities being predicted.  
	Interactions can be broken into 3D and 2D, that is whether the organism consumes from a 2D ``surface" e.g. a cow grazing or a 3D ``volume" e.g. a pelagic consumer which consumes prey from within the water column.  As such, both handling time and search rate can be defined as Eq. \ref{handling_time} and Eq. \ref{search_rate} respectively.
	\begin{equation}
		\label{search_rate}
		a(m) = a_0 m_t^\gamma
	\end{equation}
	
	\begin{equation}
		\label{handling_time}
		h(m) = t_{h,0} m_t^\beta
	\end{equation}
	A functional response alone is not enough to fully define intake.  This is because processing of consumed resources is not one hundred percent efficient which leads to inevitable loss consumed energy.  As a result, to achieve the final gain term, a dimensionless efficiency term $\epsilon$ is applied.  In poikilotherms assimilation efficiency is roughly 70\% \parencite{Peters1983}
	
	\subsubsection{Loss}
	Metabolic cost  has previously been dependant upon the gain term within traditional OGMs (see \cite{West2001, Hou2008}).  However, for non-maximal intake the relationship will no longer hold true.  As a result, this model has taken previously measured values to be used as metabolic cost (see Eq. \ref{metabolic_cost} taken from \cite{Peters1983} % add ref to hemmingsen here since that is the peters source
	and Table \ref{parameters} for further details)
	\begin{equation}
		\label{metabolic_cost}
		B_m = 0.14 m_t^\mu
	\end{equation}
	Next to take allometric scaling of reproduction into account, the reproductive cost term from \citeauthor{Charnov2001} (\citeyear{Charnov2001}) is changed from $cm^1$ which assumes isometric scaling to $cm^\rho$.  $c$ can be interpreted as the proportion of mass dedicated to reproduction, i.e. the gonadosomatic index (GSI) of the fish \parencite{Charnov2001}.  Just as in \citeauthor{Charnov2001} (\citeyear{Charnov2001}) reproductive cost is only taken into account once maturity is reached.  Meaning that until a length of time ($\alpha$) has passed, reproductive cost is considered to be zero.
	
	\subsubsection{Full Growth Equation}
	When all newly defined terms are taken together they produce Eq. \ref{dmdt_juvenile}  and \ref{dmdt_mature} which describe the change in mass for a fish before (Eq. \ref{dmdt_juvenile}) and after maturation (Eq. \ref{dmdt_mature}).
	\begin{align}
		\label{dmdt_juvenile}
		\frac{dm}{dt} &= \epsilon f(\cdot) - B_m & m_t < m_\alpha \\
		\label{dmdt_mature}
		\frac{dm}{dt} &= \epsilon f(\cdot) - B_m - cm_t^\rho & m_t \geq m_\alpha
	\end{align}

	
	
	\subsection{Reproductive Output}
	At any time ($ t $) a reproducing organism devotes some amount of energy to reproduction.  This is the product between the amount of mass dedicated to reproduction ($ cm^\rho $) and a declining efficiency term ($ h_t $) which begins at maturity ($ \alpha $) and represents reproductive senescence \parencite{Benoit2018, Vrtilek2018, Stearns2000}.  In addition to amount of reproduction, the offspring are also subject to mortality ($ l_t $).  By combining the two, lifetime reproductive output can be estimated and is described by the ``characteristic equation" (Eq. \ref{characteristic_equation}) which represents reproductive output in a non-growing population \parencite{Tsoukali2016, roff1993, Roff2001, stearns1992evolution, Arendt2011, Roff1986, Roff1984}
	\begin{equation}
		\label{characteristic_equation}
		R_0 = \int c m_t^\rho h_t l_t 
	\end{equation}
	Mortality is experienced differently by juvenile ($ t < \alpha $) and reproducing individuals ($ t \leq \alpha $). % TODO <-dont like this sentence
	Mortality of offspring prior to maturity is described as a survival rate $ l_t = e^{-Z(t)} $ which is an exponentially decreasing function bounded between zero and one.  It controls how many offspring make it to maturity.  After maturity, survival is again described as an exponential function which takes time to maturity into account, $ l_t = e^{-Z(t-\alpha)} $.  
	Reproductive senescence can be also be estimated as an exponential function which begins after maturity and declines over time  ($ e^{-k(t-\alpha)} $), where $ k $ is the senescence term.  When all values are inserted into the characteristic equation, it results in the equation used by \citeauthor{Charnov2001} (\citeyear{Charnov2001}) with the inclusion of reproductive senescence (Eq. \ref{reproductive_output}).
	\begin{equation}
		\label{reproductive_output}
		R_0 = c\int_0^\alpha e^{-Z_t} dt  \int_\alpha^\infty m_t^\rho e^{-(\kappa + Z)(t - \alpha)}dt\\
	\end{equation} 
	In Eq. \ref{reproductive_output} $ Z $ represents instantaneous mortality.  This rate has been shown to be related to time of maturation in many taxon groups, and within it follows the relationship $ \alpha \cdot Z \approx  2$.  This can then be rearranged to estimate instantaneous mortality, $ Z \approx 2/\alpha  $
	
	\subsection{Maximising Reproduction}
	It is assumed that evolution will converge on metabolic values which maximise fitness. 
	Fitness being defined as how much an individual is able to contribute to the gene pool \parencite{Speakman2008, Stearns2000}.  % Would like to remove this and link more smoothly
	To this end, lifetime reproductive output is often used as a measure of fitness \parencite{Charnov1991, Audzijonyte2018, Speakman2008, Stearns2000, Charnov2001, Tsoukali2016, Brown1993, Charnov2007}.  Therefore, by maximising for reproductive output, it should become clear what parameters will yield the highest fitness.  These parameters will then show whether, within a theoretical framework, hyperallometric scaling arises.
	
	To find all optimal values for reproduction would require Eq. \ref{reproductive_output} to analytically solved.  However, since no such solution is possible, I simulated the problem numerically to obtain a result.  This was done by simulating across values of $ c $ and $ \rho $, the parameters of interest between growth (Eq. \ref{dmdt_juvenile} and \ref{dmdt_mature}) and reproductive output (Eq. \ref{reproductive_output}).  $ c $ was bound between 0 and 0.4, which encapsulates the values measured within fish \parencite{Benoit2018, Roff1983, Fontoura2009}.  Though it has been shown to reach as much as 0.7 in some species \parencite{Parker2018}  To search for any hyperallometry within reproduction, $ \rho $ was bound between 0 and 2.  
	The simulation was then run at 0.01 value intervals in both $c$ and $\rho$ over a lifespan of one million days.  The results of each simulation were recorded and any non-viable results were discarded.  A result was considered non-viable if fish had ``shrunk" more than 5\% in order to accommodate reproductive costs.  Shrinking occurs in the model because  the combined loss of energy to metabolism and reproduction is too much for the simulated values at the mass achieved by maturation so the individual experiences a deficit of energy, which is paid by loss in mass until equilibrium is achieved. % \cite{VandenBerghe1992} for reproductive mass loss (though it is due to behaviour changes)
	Shrinking is not expected at maturity in reality.  Typically, maturity will occur while the organism still has room for growth.  It is the onset of reproduction which is considered to slow or stop growth % this is the case where the metabolic exponent is the same or less than the intake one.
	 (see Fig. \ref{OGM_Curve}).  Shrink can be thought of as starvation in a real organism.  If energetic cost are not met then energy reserves in the body, such as fat and muscle, are broken down for energy.  It has been shown that some fish are capable of losing up to 10\% of their body mass \parencite{VandenBerghe1992}.  However, this was during the breeding season and caused by behavioural changes due to parenting.  Additionally, individuals were shown to rebound back to their``normal" body mass once the breeding season had ended.
	% survival impacts of shrinking? are there sources?
	% \cite{Lambert2000}
	
	\subsection{Impact of Parameters}

	In order to determine the roles of metabolic exponent, maturation time and resource density within the model, sensitivity analyses were performed on each parameter with regard to $c$ and $\rho$.  This was done by simulating the parameters across multiple values and obtaining the optimal value for $c$ and $\rho$ as described above.
	
	
	\begin{centering}
		
		%		TODO check description capitalisation
		
		\begin{table}[h!]
			\caption{Table describing parameters used in the model, along with values, units and sources where applicable.  The units of resource density change depending on the dimension of intake.  $m^D$ represent either $m^2$ in 2D or $m^3$ in 3D} 
			\label{parameters}
			\vspace{2mm}
			\begin{tabular}{c l l l l p{3cm}}
				\hline
				Parameter 	& Description 			& Value 	& Units 	& Range 		& Source \\
				\hline
				$m$			& Mass					& -			& kg day$^{-1}$& -			&		\\
				
				$B_m$		& Metabolic Cost		& $0.14 m^{\mu}$ & kg day$^{-1}$& - 	& \cite{Peters1983}\\
				$\mu$		& Metabolic Exponent	& -			&	-		& 0.75 - 1.0	& - \\
				$\alpha$	& Age of maturity		& -     	& day		& -				& -\\
				$c$			& Reproduction scaling constant & - & kg day$^{-1}$& 0 - 0.5 		& -\\
				$\rho$		& Reproduction scaling exponent	& -	&	-		& 0 - 1.5			& -\\
				$Z$			& Rate of instantaneous mortality& $2/\alpha$	& & & \cite{Charnov2001}\\%double check ref
				$k$			& Reproductive senescence & 0.01\\
				
				$\epsilon$	& Assimilation Efficiency & 0.70 & - & - 		& \cite{Peters1983} \\
				$X_r$ 		& Resource Density		& -		& kg/m$^D$		& 0.11 - 30				& -\\
				$\gamma$	& Search rate scaling exponent & 0.68 in 2D	& - & - & \cite{Pawar2012} \\
				&						& 1.05 in 3D\\
				$a_0$		& Search rate scaling constant & $10^{-3.08}$ in 2D & m$^2$ s$^{-1}$ kg$^{-0.68}$   & - &\cite{Pawar2012}	\\
				&						& $10^{-1.77}$ in 3D& m$^2$ s$^{-1}$ kg$^{-1.05} $\\
				$\beta$		& Handling time scaling exponent& 0.75 & - & - & \cite{Pawar2012}\\
				$t_{h, 0}$	& Handling time scaling constant& $10^{3.95}$ in 2D &kg$^{1-\beta}$ s& -& \cite{Pawar2012}	\\
				&						& $10^{3.04}$ in 3D			&kg$^{1-\beta}$ s\\
				\hline
			\end{tabular}
		\end{table}
	\end{centering}
%	\subsection{notes}

	\nolinenumbers
	%%%%%%%%%% Results %%%%%%%%%%
\section{Results}
	\linenumbers
	
%	\subsection{notes}
%		scenesence not much of a factor
% 10 \% predicted under Roff 1983

	\subsection{Produced Growth Curve and Maturation Time}
	% talking about at meta_exp = 1
	As can be seen in Figure \ref{growth_curve}, growth is very fast within the model.  Asymptotic mass is reached by $\sim$15 days.  This makes interpreting any results regarding maturation time difficult % TODO < rephrase
	because any time after $\sim$15 days produces the same result (see Fig. \ref{fig:senslowresourcesmaturationtimeshortexp1} - \ref{fig:senshighresourcesmaturationtimeshortexp1}).
	However, as can be seen in Fig. \ref{fig:sensmaturationtimeshortexp075}, where maturation is occurring early during the growth phase, hyperallotric scaling emerges.  This pattern shows in particular at lower resource densities in 3D (Fig. \ref{fig:sensverylowresourcesmaturationtimeshortexp075} and \ref{fig:sensverylowresourcesmaturationtimeshortexp1})
	
	
	\subsection{Effect of Parameters}
	\subsubsection{Resource Density}
	The scaling relationship emerges as would be expected from the scaling of the functional response.  At low resource densities the output of the functional response will scale similarly to search rate.  The scaling of which is higher in 3D as resources increase and the response shifts to scaling similarly to the inverse of handling time.  At this point $\rho$ starts to take values which are higher in 2D than 3D, because of the higher normalisation constant in 2D.
	
	\subsubsection{Metabolic Exponent}
	% talking about at meta_exp = 1
	The expected result is for increasing values of metabolic scaling exponent ($\mu$) for $\rho$ to also increase.  This is because the lower values of $\mu$ will result in there being a larger gap between the scaling of intake and maintenance, which allows for steeper scaling in reproduction (see Fig. \ref{scaling_plot} b)) \parencite{Marshall2019}.  This appears to be the case when analysing $\rho$ with respect to the other parameters where $\mu = 1$ or 0.75 (e.g. Fig. \ref{fig:senshighresourcesmaturationtimeshortexp1} and \ref{fig:sensmaturationtimeshortexp075} or Fig. \ref{fig:sensresourcedensityexp1broad} and \ref{fig:sensresourcedensityexp075fine}).  However, when $\mu$ is explicitly simulated over a variety of values, the trend suggests that increasing $\mu$ allows for higher values of $\rho$.  
	
	\subsubsection{Shrinking}
	% talking about at meta_exp = 1
	At saturated resource densities, allowing for greater proportions of shrinking enables $\rho$ to be set to larger values in both dimensions.  This is perhaps not surprising as increasing the shrinking proportion allows for a larger reproductive cost since all individuals will already be at asymptotic mass at time of maturation.
	
	\subsubsection{$c$ Values}
	Estimations of $c$ are low in many cases, especially in 2D.  While this may be low compared to the $\sim$10\% - 35\% expected \parencite{Benoit2018, Fontoura2009, Roff1983} it is not unprecedented for values of 2\% to be observed in some fish \parencite{Gunderson1997}.  It may be necessary for the lower bounds of $c$ to be adjusted based on what is expected or even viable in the organisms being simulated.
	
%	minimum viable resource is dependent upon meta.exp, with 0.75 allowing for less than that of 1 and scaling in between 0.2 for 1 and 0.15 for 0.75. likewise 3D is capable of going lower than 2D due to the higher intake rate.  shrinkage plays no role here.

	\begin{figure}[h!]
				
%		TODO 1.The title for each plot needs to be cleaned up.  2.The units should be m^2 and 3.need letter labels for each.
		\includegraphics[width=\textwidth]{../../results/report_3D2D_HighLowResCLEAN_meta_exp_1.pdf}
		
		\caption{Multiplot with 2D vs 3D and varying resource densities. The optimal value is denoted by the blue circle.
		The resources used for the low resource scenario (top row) is the minimum amount of resources that allows growth with a $c$ and $\rho$ of 0.01.  As would be expected, since has steeper scaling, 3D allows for growth at smaller resource densities than 2D.
		Low reources in 2D were $ \approx 0.1kg/m^2 $ and $ 0.00035kg/m^2 $ in 3D.
		$ 100kg/m^2 $ was used for high resources because, which although unrealistic, ensures that resources are not a limiting factor in the simulations.
		Hyperallometric scaling is observed in 2D ($\rho = 1.64$ at high resources and 1.99 at low resources (NOTE: This is because 2 is the upper limit of $\rho$ I have simulated here). 
		Scaling in 3D is slightly hypoallometric $rho = 0.94$ and c = 0.01 at high resources. $\rho = 0.8$ and $c = 0.01$ at low resources
		The metabolic scaling exponent = 1 in all cases}
		\label{resources2D3D}
	\end{figure}

	\begin{figure}[h!]
		
		%		TODO 1.The title for each plot needs to be cleaned up.  2.The units should be m^2 and 3.need letter labels for each. CHECK SHRINKING
		\includegraphics[width=\textwidth]{../../results/report_3D2D_HighLowResCLEAN_meta_exp_075.pdf}
		
		\caption{Multiplot with 2D vs 3D and varying resource densities. The optimal value is denoted by the blue circle.
			The resources used for the low resource scenario (top row) is the minimum amount of resources that allows growth with a $c$ and $\rho$ of 0.01.  As would be expected, since has steeper scaling, 3D allows for growth at smaller resource densities than 2D.
			Low reources in 2D were $ \approx 0.1kg/m^2 $ and $ 0.00035kg/m^2 $ in 3D.
			$ 100kg/m^2 $ was used for high resources because, which although unrealistic, ensures that resources are not a limiting factor in the simulations.
			Hyperallometric scaling is observed in 2D ($\rho = 1.64$ at high resources and 1.99 at low resources (NOTE: This is because 2 is the upper limit of $\rho$ I have simulated here). 
			Scaling in 3D is slightly hypoallometric $rho = 0.94$ and c = 0.01 at high resources. $\rho = 0.8$ and $c = 0.01$ at low resources
			The metabolic scaling exponent = 0.75 in all cases}
		\label{resources2D3D}
	\end{figure}
	
	\begin{figure}[h!]
		
		\includegraphics[width=\textwidth]{../../results/report_2D_exponents.pdf}
		
		\caption{Multiplot with changing metabolic exponents in 2D, ranging from 0.75 to 1, resource density = 0.01, time to maturation is 5 years.  
		It can be seen that higher metabolic responses limit the proportion of mass that can be dedicated to reproduction.  However, in exchange, they allow for much steeper scaling.  A patter which is mirrored in 2D and 3D (fig. \ref{exponents3D})}
		\label{exponents2D}
	\end{figure}
		
	\begin{figure}[h!]
		
		\includegraphics[width=\textwidth]{../../results/report_3D_exponents.pdf}
		
		\caption{Multiplot with changing exponents in 3D, ranging from 0.75 to 1, resource density = 0.01, alpha is 5 years. See fig \ref{exponents2D} for explanation. (NOTE: either 2D or 3D likely belongs in SI) } %either of these likely belong in SI.
		\label{exponents3D}
	\end{figure}
		
	\begin{figure}
		\centering %TODO Units for axis?
		\includegraphics[width=\textwidth]{../../results/pretty_curve}
		\caption{Example of the growth curve and cumulative reproduction expected from a traditional OGM model. Maturation occurs at 1000 days, after which growth is less steep until reaching asymptotic mass.  }
		\label{OGM_Curve}
	\end{figure}
	
	\begin{figure}
		
		\includegraphics[width=\textwidth]{../../results/report_growth_curve.pdf}
		\caption{The growth over a fish which consumes in 2D.  Maturation occurs at 5 years (1825 days).  The fish was allowed to shrink by 5\% at the onset of reproduction.}
		\label{growth_curve}
	\end{figure}



	\nolinenumbers
	
	%%%%%%%%%% Discussion %%%%%%%%%%
\section{Discussion}
	\linenumbers
	% Overall summary of results and their meaning.
	
	
	
	% Issues with my model
		% the mass problem
			% Shrinking and how it addresses the growth speed issue in this model to a degree in my opinion.
		% gsi is likley underestimated in some cases at the very least 0.01 gsi wont be viable for all animals
	Some caveats with the results of the model.  First, growth is simulated as being extremely fast.  As previously stated, asymptotic size was reach by $\sim$10 days.  This is of course not representative of the real world, where individuals generally need several months to years to reach maturity.	
	The rapid growth may be due to several factors.  First is that metabolic cost may be underestimated.  Similar to \citeauthor{West2001} (\citeyear{West2001}), this study used resting metabolic rate to define metabolic costs.  However, this does not take other costs into account such as digestion and locomotion.
	%A rate which is typically used within energetic and growth studies, including early OGMs \parencite{Brett1971,Gillooly2001}  % ref for RMR in growth Tsoukali?  Burger? 
	This was addressed in traditional OGMs by \citeauthor{Hou2008} (\citeyear{Hou2008}).  However, due to the use of asymptotic mass in the parametrisation of this change, the same changes could not be used in this model.  Resting metabolic rate and active metabolic rates do not scale in the same way with mass \parencite{Gillooly2001, Weibel2004}.  The additional cost of active metabolic rate would cause a steeper scaling within the metabolic cost term, leading to more gradual growth (Fig. \ref{scaling_plot}). 
	As such, inclusion of active metabolic rates, while challenging to measure directly and implement, is needed.
	Additionally, there may be behavioural or physiological factors that would lead to an altered metabolic rate.  
	% TODO V presenting this as an issue but it is an area for improvement needs rephrasing and/or shifting
	In this regard, temperature plays a critical role.  It is well documented that a change in temperature will change many biological rates  (see \citeauthor{Peters1983} (\citeyear{Peters1983}), \citeauthor{Gillooly2001} (\citeyear{Gillooly2001}) etc.).  
	It has been shown that growth is dependent on temperature within fish. % sources
	For example, a 1.5\textdegree{}C increase in sea temperatures could result in a 15\% decrease in fish lengths \parencite{VanRijn2017}. % try to phrase better or bolster since it is lifted from diego
	The functional response data used in this study is standardised around 15\textdegree{}C \parencite{Pawar2012}.
	Meanwhile the metabolic cost is for an unspecified temperature \parencite{Peters1983}.  Using rates where the temperature effect is taken into account is crucial for model accuracy given the variation in rates that occurs over different temperatures.  Work such as \citeauthor{Barneche2014} (\citeyear{Barneche2014}) has investigated this effect, however the estimate for metabolic rate is several orders of magnitude lower than what was reported by \citeauthor{Peters1983} (\citeyear{Peters1983}), which is the rate used in this study.  Thus further investigation is required.  % should include a plot in the SI for this.
	
	In fish metabolic rate has also been shown to drop under starvation \parencite{Cook2000}.  In homeotherms feeding restriction has also been shown to also lower body temperature, since metabolic rate and core temperature are closely related in homeotherms \parencite{Ballor1991, Blanc2003,}.  
	
	Another possible point of error is the estimates for supply.  The parameters used from \citeauthor{Pawar2012} (\citeyear{Pawar2012}) are for a spectrum of animals from mammals to insects.  It is possible by reanalysing the data for only marine species, or more specifically only within taxon or species, predictions of supply could be improved \parencite{Marshall2019}.  % check dim paper for addressing rates in different environments
	
	A factor that is not taken into account in this model is that resources are not constant over time.  This can be implemented within the model by varying resource density over time.  The functional response will respond accordingly giving intake which varies through time.  One concern with implementing such a response is fluctuations are likely not experienced by all organisms in the same way.  For a fish with a small range a local fluctuation can be measured and described relatively simply.  However, for a fish with a very large range there is the possiblity leaving resource poor areas in search of richer waters.
	
	
	Despite these caveats, the patterns which arise from the model are promising.  The qualitative patterns seen in this study should not change even with slower growth.  This is because while the exact results are not representative of reality, the relationship between the values is.  For example, the manner in which supply and metabolic rate interact for different values of the metabolic scaling exponent.  The fact that a higher metabolic scaling exponent will cause intake and metabolic cost to intersect sooner does not change regardless of the absolute values.
	
 
	% However such field data is sparse. <- justify saying this some how. need a review or something.
	
	
	% What else can be done to address these issues
		% What is the minimum amount of GSI needed. reproductive structure will vary a bit between species and there is likely a minimum that a species must devote to these structures for them to be viable.
	
	
	% Emphasis that this is applicable to more than just fish - What needs to be done to apply the model to other animal groups.
	In conclusion, the model presented in this study is a promising base which can be expanded upon in a way that was not possible with previous OGMs allowing for much more controlled and detailed explanations of the factors controlling growth.  In contrast to all previous work, which assume optimal supply, the concept of varying supply is addressed using functional responses.  Additionally, qualitative evidence is provided supporting hyperallometric scaling in fish using energy budget as the basis.  The model can easily be applied to any animal taxon, not just fish, with some simple changes.  Additionally, there are clear directions to be explored to improve the model's accuracy.
	
	
	
	
	
	\begin{figure}
		\centering
		\includegraphics[width=\linewidth]{../../results/scalingplot.pdf}
		\caption{a) shows how maintenance cost outscales supply in a traditional OGM.  Growth only stops when maintenance reaches the supply line  b) shows scaling for supply and maintenance as equal. Since scaling is equal growth will never stop until the new cost of reproduction is introduced some time during development.  Blue line is supply, orange line is maintenance, green dotted line is reproduction}
		\label{scaling_plot}
	\end{figure}
%	\subsection{planned outline}
%	Overall summary of results and their meaning.
%	\\- I show that hyper allometry is possible within fish.  
%	\\- higher metabolic cost are leading to hyper allometry, even though I am seeing this for what the data considers too high (i.e. metabolic scaling expontent = 1) \parencite{Marshall2019} % see my google drive doc
%	\\- scaling in 3D will outscale metabolic cost in all cases Table \ref{parameters} because it is greater than 1, need reproductive cost to stop growth.
%	Shrinking and how it addresses the growth speed issue in this model to a degree in my opinion.
%	\\- Is like seeing the stage of growth we expect to see in a typical growth curve in reverse.  Where instead of growing into the asymptotic mass, it shrinks to it.  The energy balance remains the same.
%	\\
%	Issues with my model
%	- mass is too high, addressing a more detailed reproductive mass/metabolic rate may address this or reanalysing samraats data to be more taxon specific may do it.
%	\\
%	What else can be done.
%	\\- more fish specific intake using \cite{Pawar2012} methods as argued in \cite{Marshall2019}
%	\\- trying to model specific species with know growth curves to see how well our results map to them even if the curve isn't great it could be that reproductive output and final mass map well.
%	\\- more detailed metabolic cost since here is basically just RMR with higher exponents.  It may be that the constant shift too and this would lead to better predictions
%	\\- Incorporate oscillating resource density.  was attempted here but due to computational problems could not be implemented.  There is a lack of oscillation data this though and the goal would need to be far more specific since it will depend on scale and range of the organism.
%	\\- This same framework could be applied to other taxa with just a little adjustment (metabolic cost and reproductive equation?)
%	\\- There may be behavioural changes that occur at under feeding restriction that slows metabolism, similar to hibernation etc.
%	\\- temperature could be added here to adjust the metabolic and functional response rates %samraats data is standardised to 15C doesn't seem to be noted in peter's, hemingsen may show it
%	- develop the model to include storage like DEB or put this into a DEB framework.
%	\\
%	What is the minimum amount of GSI needed. reproductive structure will vary a bit between species and there is likely a minimum that a species must devote to these structures for them to be viable. %\cite{Gunderson1997} shows that mortality increases with higher GSI and larger inds have lower GSI does show low values for GSI (0.02) 
%	%GSI can reach 0.75 in invertebrates \cite{Parker2018}
%	\\
%	Emphasis that this is applicable to more than just fish - What needs to be done to apply the model to other animal groups. % need to open up gsi \cite{Parker2018} 0.75 inverts
%	\\
%	This is all energetically based, allowing us to ask questions of why or how there are exceptions to this in the case that organsims should be able to grow larger energetically speaking but may be limited genetically or via some other mechanism
	
	\nolinenumbers
%	\subsection{notes}
%	there are some integration issues that could allow for the appearance of more possible values
%	The shrinking helps to show what could be possible given there is a more gentle growth curve since the main difference is that the time when growth become more limited is just shown as shrinking vs slower growth.


	%%%%%%%%%% Conclusion %%%%%%%%%%
\section{Conclusion}
	\linenumbers
	
	
	
	\nolinenumbers
	
	%%%%%%%%%% Data Declaration %%%%%%%%%%
	\textbf{Code and Data Availability}

	Code is available at: \url{https://github.com/Don-Burns/Masters_Project}
	
	
	%%%%%%%%%% Bibliography %%%%%%%%%%
	\newpage
	\linenumbers
	
	%nicer indentation
	\addcontentsline{toc}{section}{\protect\numberline{}References}	
	
	\printbibliography
	
	% prob more official method of doing the above
%	\printbibliography[heading=bibintoc, title={References}]
	\nolinenumbers
	
	
	%%%%%%%%%% SI %%%%%%%%%%
	\newpage
	\include{SI}
	
	
\end{document}} % to include word count 

%%%%%%%%%% Formatting %%%%%%%%%%
\usepackage[margin=2cm]{geometry} % margins of 2cm
\linespread{1.5} %1.5 spacing
%\renewcommand{\familydefault}{\sfdefault} % set font to arial clone (helvet)

\usepackage[compact]{titlesec} % reduce spacing bewteen section titles

%%%%%%%%%% Bibliography %%%%%%%%%%
\addbibresource{../Masters_Thesis.bib}

%commands 

%%%%%%%%%% Document %%%%%%%%%%
\begin{document}
	%%%%%%%%%% Title Page %%%%%%%%%%
	%\newcommand{\crest}{\includegraphics[width = 4cm, keepaspectratio]{../images/IC_Crest.eps}} % Imperial crest
%%formating
\begin{titlepage} % Suppresses headers and footers on the title page
	
	\centering % Centre everything on the title page
	
	\scshape % Use small caps for all text on the title page
	
%	\vspace*{\baselineskip} % White space at the top of the page
	
	%------------------------------------------------
	%	Title
	%------------------------------------------------
	
	\rule{\textwidth}{1.6pt}\vspace*{-\baselineskip}\vspace*{2pt} % Thick horizontal rule
	\rule{\textwidth}{0.4pt} % Thin horizontal rule
	
	\vspace{0.75\baselineskip} % Whitespace above the title
	
	{\LARGE Using Realistic Resource Supply to Show Hyperallometric Reproductive Output in Fish\\} % Title
	
	\vspace{0.75\baselineskip} % Whitespace below the title
	
	\rule{\textwidth}{0.4pt}\vspace*{-\baselineskip}\vspace{3.2pt} % Thin horizontal rule
	\rule{\textwidth}{1.6pt} % Thick horizontal rule
	
	\vspace{1\baselineskip} % Whitespace after the title block
	
	%------------------------------------------------
	%	Subtitle
	%------------------------------------------------
	
	%SUBTITLE? % Subtitle or further description
%	Student:
	
	
	\vspace{0.5\baselineskip} % Whitespace before 
	
	{\scshape\Large D\'onal Burns  \\} % my name
	
	\vspace{0.5\baselineskip} % Whitespace below 
	
	\textit{CID: 01749638 \\ Imperial College London \\ Email: donal.burns@imperial.ac.uk} % affiliation and email
	
	\vspace*{2\baselineskip} % Whitespace under the subtitle
	
	
	
%	Supervisor:
%	
%	
%	\vspace{0.5\baselineskip} % Whitespace before 
%	
%	{\scshape\Large Samraat Pawar \\} % supervisor name
%	
%	\vspace{0.5\baselineskip} % Whitespace below 
%	
%	\textit{Imperial College London \\ Email: s.pawar@imperial.ac.uk} % affiliation and email
%	
%	\vspace{3cm} % Whitespace between 
	

	
	%% crest
	
	\includegraphics[width = 4cm, keepaspectratio]{../images/IC_crest.pdf}
	
	\vspace{0.3\baselineskip} % Whitespace under the Uni logo
	
	Submitted: August 27$^{th}$ 2020 % Publication Date
	
	

		%%Submission clause
	\vspace{0.5\baselineskip}
	
	A thesis submitted in partial fulfilment of the requirements for the degree of
	%Computational Methods in Ecology and Evolution 
	Master of Science at Imperial College London
	\vspace{0.5\baselineskip}
	
	Formatted in the journal style of Functional Ecology	
	\vspace{0.5\baselineskip}
	
	Submitted for the MSc in Computational Methods in Ecology and Evolution
	
\end{titlepage}

	

	%%%%%%%%%% Declaration %%%%%%%%%%
	\section*{Declaration}
	I declare this project as my own work.  The model presented here was developed in conjunction with my supervisor, Dr. Samraat Pawar, and Ph.D. students Tom Clegg and Olivia Morris.  I was responsible for any simulations and data presentation.\newline
	%% word count
	\textbf{Word Count: \wordcount}

	\newpage
	
	%%%%%%%%%% Abstract %%%%%%%%%%
	\section*{Abstract}
	\linenumbers
%	Size is essential to reproductive output. By extension understanding growth, determines size allows understanding of reproductive output. 
	 
	With recent results showing that reproduction in fish scales hyperallometrically there is a need to update growth OGMs to reflect this fact.  Current OGMs assume optimal intake, an assumption which is not always reflected in the field.  This study develops an energy intake focused approach to explaining growth, an area which has not been covered within current literature, and shows that hyperallometric scaling of reproductive output arises when allowing for variable reproductive scaling and maximising for fitness.  The model is applicable to not only fish, but any animals taxon with some simple parameter adjustments.  I offer direction for improvements and areas to be developed in order to allow the model to be applicable to any temperature range.
	\vspace*{0.5 cm}
	\newline
%	TODO pieces to include:
%	
%	show the effects of the exponent and why it is important that we use the correct one moving forward.
%	give direction as to what needs to be done going forward.  what areas need work
	%TODO KEywords
	\textbf{Keywords:}\\
	allometry; functional response; metabolic theory; growth; intake; life history; metabolism; reproduction; reproductive output; supply

	
	
	\nolinenumbers
	%%%%%%%%%% Acknowledgements %%%%%%%%%%
	%\thispagestyle{empty}

\mbox{}\newline\vspace{10mm} \mbox{}\LARGE
%
{\bf Acknowledgements} \normalsize \vspace{5mm}

%I would like to thank my supervisor.....
%thanks to luke who gave a great base to start from on the project





	
	%%%%%%%%%% Table of Contents %%%%%%%%%%
	\tableofcontents
	\newpage

	%%%%%%%%%% Introduction %%%%%%%%%%
\section{Introduction}
	\linenumbers
	% putting OG first
	Growth is a gateway to understanding how many biological rates change throughout ontogeny.  This is because so many biological rates scale with mass \parencite{Kleiber1932}.  To date many models have been developed to predict and describe the growth of an organism throughout its lifetime.  The three main approaches used are the von Bertalanffy model, the dynamic energy budget (DEB) model and the ontogenetic growth model (OGM), which will be the focus of this study \parencite{Putter1918, vonBertalanffy1938, Kooijman1986, West2001}.  All of these are energetic based models with varying assumptions, key among which is the scaling of resource supply and metabolic rate with mass.   In OGMs supply is thought of as being optimal at all times which leads to the assumption that intake scales to the power of 0.75.  Indeed while under optimal conditions this may be true, however, it neglects that this situation is thought to rarely occur in the field \parencite{Pawar2012}.
%	Additionally, growth model approaches have made various assumptions about the scaling of metabolic rates.  OGMs assume that scaling is isometric with mass, i.e. scaling exponent is 1.  This is contrary to what has been thought 
	This study takes the novel approach of using a mass specific functional response and assimilation efficiency to describe how intake changes both throughout ontogeny and varying levels of resource availability. While the methodology used here can be used for any animal taxon, in order to investigate growth, a taxon much be chosen to parametrise metabolic rate around.  This study will focus on fish.
		
% TODO bring in OGM growth curve and use to illustate the basic OGM as per samraats comments
	% Why is growth important
	Understanding how organisms grow and what factors play a role in determining growth is not only insightful from the perspective of understanding the world around us, but can also be used to better manage the many fisheries and marine protected areas around the world.  An objective which is becoming increasingly important as the oceans' fish stocks continue to be depleted by over fishing. 
	%VV rephrase this part to hint more towards size than it currently does VV
	This is compounded by global warming which threatens to alter the structure of the ecosystems even if they are not fished and left in their ``natural" state. 
	
	Larger fish produce more offspring than smaller ones.  % this is known in general isometrically
	In fact, it has been shown that larger fish produce far more offspring than the equivalent mass composed of smaller fish,  that is a 2kg fish will produce more offspring than two 1kg fish \parencite{Barneche2018}.
	Larger fish also use energy more efficiently than multiple smaller ones per unit mass.  This is due to them having a lower mass specific metabolic rate \parencite{Peters1983, Kleiber1932, Brown2004}.  
	%As a result a higher proportion of larger fish may actually reduce stress upon the ecosystem from an energetic perspective which could be useful when trying to manage fish stocks. %need citations for this also questionably neccessary as a point and has little backup? may even not be the case at all due to other pressures introduced by larger individuals in a system
	% larger mothers produce larger offspring which may better survive \cite{Barneche2018}
	Additionally, larger mothers produce larger offspring, which are then more likely to survive to adulthood and reproduce \parencite{Hixon2014, Marshall2006}. % double check this statement from barneche and get more backing
	%
	
	It is already known that metabolic rate is dependant on temperature which in turn will lower fish sizes \parencite{Gillooly2001, Brown2004} and with increasing global temperatures, understanding in greater detail how increased metabolic rates may affect growth is useful in population management.
	%
	
%	Growth is in essence a balancing act between how much energy an organism can acquire and the amount of energy required for maintenance, 
%	that is movement, digestion et cetera.  
	
	%power laws -- needs citation or could be counted as common knowledge? \cite{Huxley1932}
	Key to understanding biological rates and their relationship with mass is the concept of power laws. Many biological traits can be described as scaling to the power of some other biological trait.  That is some rate, $Y$, can be expressed for any mass by $Y = Y_0 m^\beta$, where $m$ is mass \parencite{Kleiber1932}.  Power laws can be broadly categorised based on the value of their exponent, $\beta$.  
	Where the exponent equals one the relationship is described as isometric or linear, that is the two traits increase at the same rate.  
	Where the exponent does not equal one, the relationship is said to be allometric.  That is the trait does not increase at the same rate as the trait being compared against.  In cases where the exponent is greater than one, the trait is said to scale super-linearly or hyperallometrically and to scale sub-linearly or hypoallometrically for cases with an exponent of less than one.
	For the purposes of describing growth these relationships are indispensable.   
	
	
%	Traditionally ontogenetic growth models have relied on knowing how large an organism is expected to grow.
%	%	von Bertalanffy / puetter model
%	The von Bertalanffy growth equation relies on knowing the longest a fish can be and the length of the fish at the beginning of measurements \parencite{vonBertalanffy1938, Putter1918}.  From here, for a known growth rate, the length of the fish after a set amount of time has passed can be predicted.
%	
%	%OGMs
%	Another approach has been to try explain growth from an energetic perspective.  That is an individuals growth is determined by the amount of energy available to it.  Once all costs, such as maintenance, have been paid all remaining energy can be invested towards growth.  There are two frameworks which have been built around this concept; dynamic energy budget (DEB) theory % TODO get ref for DEB original/important sources %add note to see discussion for more on DEB?
%	and ontogenetic growth models (OGMs), which will be focused on here.

	%	introduce \cite{West2001}	
	One of the best known examples of an OGM is the model developed by \cite{West2001}.  This model is parametrised around the average energy content of animal tissue and asymptotic mass.  Asymptotic mass being the mass at which growth has essentially stopped due to metabolic cost and energy intake equalling each other (Fig. \ref{scaling_plot}a). The model hinges on the scaling with mass between energy intake (m$^{0.75}$, i.e. allometric sub-linear scaling) and maintenance cost (m$^1$, i.e. isometric linear scaling).  In other words, as mass increases maintenance costs will slowly overtake the intake rate and halt growth.  	
	
	% talk about determinant and indeterminant growth/
	% move to including reproduction with \cite{Charnov2001} mashed with \cite{Hou2008} imporvements briefly
	The framework used by \citeauthor{West2001} (\citeyear{West2001}) was latter developed by \citeauthor{Charnov2001} (\citeyear{Charnov2001}) to take the cost of reproduction into account and allowing the estimation of lifetime production of offspring.  \citeauthor{Hou2008} (\citeyear{Hou2008}) latter developed \citeauthor{West2001}'s model further by expanding maintenance cost to include the cost of feeding and digestion (specific dynamic action), synthesis of new tissue and activity.
	
%	% move to including reproduction with \cite{Charnov2001}
%	The framework used by \citeauthor{West2001} (\citeyear{West2001}) was latter developed by \citeauthor{Charnov2001} (\citeyear{Charnov2001}) to take the cost of reproduction into account and allowing the estimation of lifetime production of offspring.  This was done by adding an additional cost which dedicated some portion of mass to reproductive output, essentially gonadosomatic index.	%like 20% max? in fish very high in  invertabrates \cite{Parker2018}
%	% improved by hou et al \cite{Hou2008} talk about including SDA and energy improvements
%	\citeauthor{West2001}'s (\citeyear{West2001}) model was later expanded upon by \citeauthor{Hou2008} \citeyear{Hou2008} to elaborate on the maintenance cost.  \citeauthor{Hou2008} (\citeyear{Hou2008}) takes both the difference between the cost of maintaining tissue and creating it, a factor that is neglected in \citeauthor{West2001} (\citeyear{West2001}), and the cost of digestion into account.
	
	% begin caveats
	% tautology still present 

	% Discuss allometry and isometry here to highlight what scaling super or sub linearly means
	In the above OGMs intake is assumed to scale sub-linearly to the power of 0.75.  This is due to the assumption that the individuals are consuming at an optimal rate at all times and therefore the only limitation is their ability to make use of that energy.  In this case, intake should theoretically scale to the power of 0.75 (see \cite{West1997}).  However, this is not always the case in the field.  It has been shown that, for non-optimal consumption, steeper scaling can occur \parencite{Pawar2012, Peters1983}. Additionally, OGMs, like many growth and metabolic models, typically use basal or resting metabolic rate to calculate metabolic cost.  This is the minimal metabolic rate of an organism and is typically thought of as the rate of the organism when at relaxed and at rest.  However, it has been shown, once factors such as movement are taken into account, that the scaling becomes steeper \parencite{Weibel2004}.
	% still no accounting for imperfect intake though is somewhat covered by \cite{Hou2011}
	% bring up the caveats/assumptions of the model i.e. isometric metabolism and repro, 0.75 scaling of intake. optimal intake at all times etc. 
	% tautology of OGMs and \cite{Hou2011} close but still issues
	The issue of non-optimal feeding is addressed somewhat by \citeauthor{Hou2011} (\citeyear{Hou2011}).  However, this growth was only investigated as, essentially, a proportion of optimal consumption and does not address a potential change in scaling of intake rate.
	Another limitation of the models used in previous OGMs is the dependence on asymptotic mass.  The models are entirely dependant on the value of optimal intake and asymptotic mass.  All other values, such as metabolic cost, are then derived in relation to these.  However, organisms are not born with some inherent restriction to the size they can attain, at least not energetically.  If there is surplus energy for a given mass the organism should be able to grow.  Relying on asymptotic mass to define the upper bound of attainable mass does not allow for investigation of the mechanisms that underpin asymptotic size in reality. 
	
	% talk about how we assume that evolution settles upon values which optimise for reproduction -- can do this when talking about charnov?
	
	% Justify and appeal my methods	
	With two of the key assumptions of current OGMs, that reproduction and metabolism scale isometrically, not holding in the field \parencite{Barneche2018, Pawar2012, Peters1983} there is a need to take an unexplored approach to modelling fish growth, in particular choosing to focus on developing how intake is described so as to better reflect the real world.  To do this an obvious starting point is to model intake as a functional response \parencite{Holling1959} so as to better reflect real world intake rates in terms of consumed biomass over time.  Non-optimal supply is a currently unexplored area within growth modelling.  This is likely due to the difficulty of directly measuring intake, especially in the field, Perhaps as a result comparatively less is known about consumption.  This leads to the need for proxy values to estimate it, for example nutrient flux \parencite{Schiettekatte2020}, or drawing broad relationships to approximate consumption, as this study will do.
	Changing the manner in which intake is defined also requires changing metabolic cost, because the two are dependent upon each other in current OGMs.  This can be done by defining metabolic rate as a value dependent on current mass rather than asymptotic mass as has been done in OGMs up until this point.  This thought process is more mechanistic since an organism has no concept of ``How large should I grow?", but rather will acquire as much resources as it capable of at its current life stage and size.  Taking this more bottom-up mechanistic approach also allows exploration of what controls how large fish grow to, since as previously mentioned, from an energetic standpoint, an organism can grow indefinitely as long as there is surplus energy available after costs have been paid.  Of course, there are also mechanical and genetic limitations upon organism size, but once size is constrained to what is known to exist, this should not be an issue.  
	Under the assumption that fish have evolved to maximise reproductive output, using the above framework, simulations can be carried out to demonstrate what conditions need to be met in order to achieve hyperallometric scaling of reproduction from an energetic perspective.  This study will answer 1) Under what metabolic scaling exponent can hyperallometric reproduction arise, 2) 
	% describe basic OGM and how i change the intake term to reflect supply
	
%	\cite{Sibly2020} % for metabolic scaling
%	\cite{Marshall2019} non-isometric scaling of repro#

% questions
% what effect does exponent actually have in relation to growth
% what coditions need to be fuflilled to possibly acheive hyperallometry
% 
	\nolinenumbers
	
	%%%%%%%%%% Methods %%%%%%%%%%
\section{Methods}
	\linenumbers
	
	\subsection{Altering OGMs to account for supply}
	In order to address the issue of supply in the context of an OGM, which can be generically described as $dm/dt = gain - loss$, some changes need to be made to the model's terms.  The first is to remove the assumption of asymptotic mass and the reliance of cost upon it.  Within a tradional OGM the gain term ($ a $) and asymptotic mass are used to define the metabolic cost ($ b $).  However, since the assumption of perfect intake is going to be broken, because of the variable supply, this relationship no longer holds.  As such, both intake and metabolic cost need to be redefined.  Additionally, in light of recent work showing that reproduction scales allometrically and not isometrically, the reproductive cost must also be modified from the form used by \citeauthor{Charnov2001} (\citeyear{Charnov2001}) \parencite{Marshall2019, Barneche2018} %TODO check marshal ref data origin
	
	
	\subsubsection{Gain}
	To define intake a natural starting place is the functional response \parencite{Holling1959}.  Functional responses  are used to define how much an organism consumes for a given resource density and is described by the following equation:	
	\begin{equation}
		\label{functional_repsonse}
		f(\cdot) = \frac{a X_r}{1 + a h X_r}
	\end{equation}
	where, $ f(\cdot) $ is the functional response, $ a $ is the search rate, $ h $ is handling time and $ X_r $ is resource density.  
	For a fixed mass and increasing resource density Eq \ref{functional_repsonse} produces a sigmoidal shape with intake eventually reaching an asymptote after some saturating amount of resources is reached.  At lower resource densities, the intake rate is primarily defined by the search rate with higher search rates yielding higher intake rates.  Conversely, at high resource densities, intake rate is approximately equal to the inverse of the handling time ($ h^{-1} $), where lower handling times yield higher intake rates.  
	
	An organism's functional response will not remain constant throughout its life history.  Search rate and handling time are affected by both the organism's mass and how it interacts with its environment \parencite{Pawar2012}.  
%	Within this model mass will be known for all time points since that is one of the quantities being predicted.  
	Interactions can be broken into 3D and 2D, that is whether the organism consumes from a 2D ``surface" e.g. a cow grazing or a 3D ``volume" e.g. a pelagic consumer which consumes prey from within the water column.  As such, both handling time and search rate can be defined as Eq. \ref{handling_time} and Eq. \ref{search_rate} respectively.
	\begin{equation}
		\label{search_rate}
		a(m) = a_0 m_t^\gamma
	\end{equation}
	
	\begin{equation}
		\label{handling_time}
		h(m) = t_{h,0} m_t^\beta
	\end{equation}
	A functional response alone is not enough to fully define intake.  This is because processing of consumed resources is not one hundred percent efficient which leads to inevitable loss consumed energy.  As a result, to achieve the final gain term, a dimensionless efficiency term $\epsilon$ is applied.  In poikilotherms assimilation efficiency is roughly 70\% \parencite{Peters1983}
	
	\subsubsection{Loss}
	Metabolic cost  has previously been dependant upon the gain term within traditional OGMs (see \cite{West2001, Hou2008}).  However, for non-maximal intake the relationship will no longer hold true.  As a result, this model has taken previously measured values to be used as metabolic cost (see Eq. \ref{metabolic_cost} taken from \cite{Peters1983} % add ref to hemmingsen here since that is the peters source
	and Table \ref{parameters} for further details)
	\begin{equation}
		\label{metabolic_cost}
		B_m = 0.14 m_t^\mu
	\end{equation}
	Next to take allometric scaling of reproduction into account, the reproductive cost term from \citeauthor{Charnov2001} (\citeyear{Charnov2001}) is changed from $cm^1$ which assumes isometric scaling to $cm^\rho$.  $c$ can be interpreted as the proportion of mass dedicated to reproduction, i.e. the gonadosomatic index (GSI) of the fish \parencite{Charnov2001}.  Just as in \citeauthor{Charnov2001} (\citeyear{Charnov2001}) reproductive cost is only taken into account once maturity is reached.  Meaning that until a length of time ($\alpha$) has passed, reproductive cost is considered to be zero.
	
	\subsubsection{Full Growth Equation}
	When all newly defined terms are taken together they produce Eq. \ref{dmdt_juvenile}  and \ref{dmdt_mature} which describe the change in mass for a fish before (Eq. \ref{dmdt_juvenile}) and after maturation (Eq. \ref{dmdt_mature}).
	\begin{align}
		\label{dmdt_juvenile}
		\frac{dm}{dt} &= \epsilon f(\cdot) - B_m & m_t < m_\alpha \\
		\label{dmdt_mature}
		\frac{dm}{dt} &= \epsilon f(\cdot) - B_m - cm_t^\rho & m_t \geq m_\alpha
	\end{align}

	
	
	\subsection{Reproductive Output}
	At any time ($ t $) a reproducing organism devotes some amount of energy to reproduction.  This is the product between the amount of mass dedicated to reproduction ($ cm^\rho $) and a declining efficiency term ($ h_t $) which begins at maturity ($ \alpha $) and represents reproductive senescence \parencite{Benoit2018, Vrtilek2018, Stearns2000}.  In addition to amount of reproduction, the offspring are also subject to mortality ($ l_t $).  By combining the two, lifetime reproductive output can be estimated and is described by the ``characteristic equation" (Eq. \ref{characteristic_equation}) which represents reproductive output in a non-growing population \parencite{Tsoukali2016, roff1993, Roff2001, stearns1992evolution, Arendt2011, Roff1986, Roff1984}
	\begin{equation}
		\label{characteristic_equation}
		R_0 = \int c m_t^\rho h_t l_t 
	\end{equation}
	Mortality is experienced differently by juvenile ($ t < \alpha $) and reproducing individuals ($ t \leq \alpha $). % TODO <-dont like this sentence
	Mortality of offspring prior to maturity is described as a survival rate $ l_t = e^{-Z(t)} $ which is an exponentially decreasing function bounded between zero and one.  It controls how many offspring make it to maturity.  After maturity, survival is again described as an exponential function which takes time to maturity into account, $ l_t = e^{-Z(t-\alpha)} $.  
	Reproductive senescence can be also be estimated as an exponential function which begins after maturity and declines over time  ($ e^{-k(t-\alpha)} $), where $ k $ is the senescence term.  When all values are inserted into the characteristic equation, it results in the equation used by \citeauthor{Charnov2001} (\citeyear{Charnov2001}) with the inclusion of reproductive senescence (Eq. \ref{reproductive_output}).
	\begin{equation}
		\label{reproductive_output}
		R_0 = c\int_0^\alpha e^{-Z_t} dt  \int_\alpha^\infty m_t^\rho e^{-(\kappa + Z)(t - \alpha)}dt\\
	\end{equation} 
	In Eq. \ref{reproductive_output} $ Z $ represents instantaneous mortality.  This rate has been shown to be related to time of maturation in many taxon groups, and within it follows the relationship $ \alpha \cdot Z \approx  2$.  This can then be rearranged to estimate instantaneous mortality, $ Z \approx 2/\alpha  $
	
	\subsection{Maximising Reproduction}
	It is assumed that evolution will converge on metabolic values which maximise fitness. 
	Fitness being defined as how much an individual is able to contribute to the gene pool \parencite{Speakman2008, Stearns2000}.  % Would like to remove this and link more smoothly
	To this end, lifetime reproductive output is often used as a measure of fitness \parencite{Charnov1991, Audzijonyte2018, Speakman2008, Stearns2000, Charnov2001, Tsoukali2016, Brown1993, Charnov2007}.  Therefore, by maximising for reproductive output, it should become clear what parameters will yield the highest fitness.  These parameters will then show whether, within a theoretical framework, hyperallometric scaling arises.
	
	To find all optimal values for reproduction would require Eq. \ref{reproductive_output} to analytically solved.  However, since no such solution is possible, I simulated the problem numerically to obtain a result.  This was done by simulating across values of $ c $ and $ \rho $, the parameters of interest between growth (Eq. \ref{dmdt_juvenile} and \ref{dmdt_mature}) and reproductive output (Eq. \ref{reproductive_output}).  $ c $ was bound between 0 and 0.4, which encapsulates the values measured within fish \parencite{Benoit2018, Roff1983, Fontoura2009}.  Though it has been shown to reach as much as 0.7 in some species \parencite{Parker2018}  To search for any hyperallometry within reproduction, $ \rho $ was bound between 0 and 2.  
	The simulation was then run at 0.01 value intervals in both $c$ and $\rho$ over a lifespan of one million days.  The results of each simulation were recorded and any non-viable results were discarded.  A result was considered non-viable if fish had ``shrunk" more than 5\% in order to accommodate reproductive costs.  Shrinking occurs in the model because  the combined loss of energy to metabolism and reproduction is too much for the simulated values at the mass achieved by maturation so the individual experiences a deficit of energy, which is paid by loss in mass until equilibrium is achieved. % \cite{VandenBerghe1992} for reproductive mass loss (though it is due to behaviour changes)
	Shrinking is not expected at maturity in reality.  Typically, maturity will occur while the organism still has room for growth.  It is the onset of reproduction which is considered to slow or stop growth % this is the case where the metabolic exponent is the same or less than the intake one.
	 (see Fig. \ref{OGM_Curve}).  Shrink can be thought of as starvation in a real organism.  If energetic cost are not met then energy reserves in the body, such as fat and muscle, are broken down for energy.  It has been shown that some fish are capable of losing up to 10\% of their body mass \parencite{VandenBerghe1992}.  However, this was during the breeding season and caused by behavioural changes due to parenting.  Additionally, individuals were shown to rebound back to their``normal" body mass once the breeding season had ended.
	% survival impacts of shrinking? are there sources?
	% \cite{Lambert2000}
	
	\subsection{Impact of Parameters}

	In order to determine the roles of metabolic exponent, maturation time and resource density within the model, sensitivity analyses were performed on each parameter with regard to $c$ and $\rho$.  This was done by simulating the parameters across multiple values and obtaining the optimal value for $c$ and $\rho$ as described above.
	
	
	\begin{centering}
		
		%		TODO check description capitalisation
		
		\begin{table}[h!]
			\caption{Table describing parameters used in the model, along with values, units and sources where applicable.  The units of resource density change depending on the dimension of intake.  $m^D$ represent either $m^2$ in 2D or $m^3$ in 3D} 
			\label{parameters}
			\vspace{2mm}
			\begin{tabular}{c l l l l p{3cm}}
				\hline
				Parameter 	& Description 			& Value 	& Units 	& Range 		& Source \\
				\hline
				$m$			& Mass					& -			& kg day$^{-1}$& -			&		\\
				
				$B_m$		& Metabolic Cost		& $0.14 m^{\mu}$ & kg day$^{-1}$& - 	& \cite{Peters1983}\\
				$\mu$		& Metabolic Exponent	& -			&	-		& 0.75 - 1.0	& - \\
				$\alpha$	& Age of maturity		& -     	& day		& -				& -\\
				$c$			& Reproduction scaling constant & - & kg day$^{-1}$& 0 - 0.5 		& -\\
				$\rho$		& Reproduction scaling exponent	& -	&	-		& 0 - 1.5			& -\\
				$Z$			& Rate of instantaneous mortality& $2/\alpha$	& & & \cite{Charnov2001}\\%double check ref
				$k$			& Reproductive senescence & 0.01\\
				
				$\epsilon$	& Assimilation Efficiency & 0.70 & - & - 		& \cite{Peters1983} \\
				$X_r$ 		& Resource Density		& -		& kg/m$^D$		& 0.11 - 30				& -\\
				$\gamma$	& Search rate scaling exponent & 0.68 in 2D	& - & - & \cite{Pawar2012} \\
				&						& 1.05 in 3D\\
				$a_0$		& Search rate scaling constant & $10^{-3.08}$ in 2D & m$^2$ s$^{-1}$ kg$^{-0.68}$   & - &\cite{Pawar2012}	\\
				&						& $10^{-1.77}$ in 3D& m$^2$ s$^{-1}$ kg$^{-1.05} $\\
				$\beta$		& Handling time scaling exponent& 0.75 & - & - & \cite{Pawar2012}\\
				$t_{h, 0}$	& Handling time scaling constant& $10^{3.95}$ in 2D &kg$^{1-\beta}$ s& -& \cite{Pawar2012}	\\
				&						& $10^{3.04}$ in 3D			&kg$^{1-\beta}$ s\\
				\hline
			\end{tabular}
		\end{table}
	\end{centering}
%	\subsection{notes}

	\nolinenumbers
	%%%%%%%%%% Results %%%%%%%%%%
\section{Results}
	\linenumbers
	
%	\subsection{notes}
%		scenesence not much of a factor
% 10 \% predicted under Roff 1983

	\subsection{Produced Growth Curve and Maturation Time}
	% talking about at meta_exp = 1
	As can be seen in Figure \ref{growth_curve}, growth is very fast within the model.  Asymptotic mass is reached by $\sim$15 days.  This makes interpreting any results regarding maturation time difficult % TODO < rephrase
	because any time after $\sim$15 days produces the same result (see Fig. \ref{fig:senslowresourcesmaturationtimeshortexp1} - \ref{fig:senshighresourcesmaturationtimeshortexp1}).
	However, as can be seen in Fig. \ref{fig:sensmaturationtimeshortexp075}, where maturation is occurring early during the growth phase, hyperallotric scaling emerges.  This pattern shows in particular at lower resource densities in 3D (Fig. \ref{fig:sensverylowresourcesmaturationtimeshortexp075} and \ref{fig:sensverylowresourcesmaturationtimeshortexp1})
	
	
	\subsection{Effect of Parameters}
	\subsubsection{Resource Density}
	The scaling relationship emerges as would be expected from the scaling of the functional response.  At low resource densities the output of the functional response will scale similarly to search rate.  The scaling of which is higher in 3D as resources increase and the response shifts to scaling similarly to the inverse of handling time.  At this point $\rho$ starts to take values which are higher in 2D than 3D, because of the higher normalisation constant in 2D.
	
	\subsubsection{Metabolic Exponent}
	% talking about at meta_exp = 1
	The expected result is for increasing values of metabolic scaling exponent ($\mu$) for $\rho$ to also increase.  This is because the lower values of $\mu$ will result in there being a larger gap between the scaling of intake and maintenance, which allows for steeper scaling in reproduction (see Fig. \ref{scaling_plot} b)) \parencite{Marshall2019}.  This appears to be the case when analysing $\rho$ with respect to the other parameters where $\mu = 1$ or 0.75 (e.g. Fig. \ref{fig:senshighresourcesmaturationtimeshortexp1} and \ref{fig:sensmaturationtimeshortexp075} or Fig. \ref{fig:sensresourcedensityexp1broad} and \ref{fig:sensresourcedensityexp075fine}).  However, when $\mu$ is explicitly simulated over a variety of values, the trend suggests that increasing $\mu$ allows for higher values of $\rho$.  
	
	\subsubsection{Shrinking}
	% talking about at meta_exp = 1
	At saturated resource densities, allowing for greater proportions of shrinking enables $\rho$ to be set to larger values in both dimensions.  This is perhaps not surprising as increasing the shrinking proportion allows for a larger reproductive cost since all individuals will already be at asymptotic mass at time of maturation.
	
	\subsubsection{$c$ Values}
	Estimations of $c$ are low in many cases, especially in 2D.  While this may be low compared to the $\sim$10\% - 35\% expected \parencite{Benoit2018, Fontoura2009, Roff1983} it is not unprecedented for values of 2\% to be observed in some fish \parencite{Gunderson1997}.  It may be necessary for the lower bounds of $c$ to be adjusted based on what is expected or even viable in the organisms being simulated.
	
%	minimum viable resource is dependent upon meta.exp, with 0.75 allowing for less than that of 1 and scaling in between 0.2 for 1 and 0.15 for 0.75. likewise 3D is capable of going lower than 2D due to the higher intake rate.  shrinkage plays no role here.

	\begin{figure}[h!]
				
%		TODO 1.The title for each plot needs to be cleaned up.  2.The units should be m^2 and 3.need letter labels for each.
		\includegraphics[width=\textwidth]{../../results/report_3D2D_HighLowResCLEAN_meta_exp_1.pdf}
		
		\caption{Multiplot with 2D vs 3D and varying resource densities. The optimal value is denoted by the blue circle.
		The resources used for the low resource scenario (top row) is the minimum amount of resources that allows growth with a $c$ and $\rho$ of 0.01.  As would be expected, since has steeper scaling, 3D allows for growth at smaller resource densities than 2D.
		Low reources in 2D were $ \approx 0.1kg/m^2 $ and $ 0.00035kg/m^2 $ in 3D.
		$ 100kg/m^2 $ was used for high resources because, which although unrealistic, ensures that resources are not a limiting factor in the simulations.
		Hyperallometric scaling is observed in 2D ($\rho = 1.64$ at high resources and 1.99 at low resources (NOTE: This is because 2 is the upper limit of $\rho$ I have simulated here). 
		Scaling in 3D is slightly hypoallometric $rho = 0.94$ and c = 0.01 at high resources. $\rho = 0.8$ and $c = 0.01$ at low resources
		The metabolic scaling exponent = 1 in all cases}
		\label{resources2D3D}
	\end{figure}

	\begin{figure}[h!]
		
		%		TODO 1.The title for each plot needs to be cleaned up.  2.The units should be m^2 and 3.need letter labels for each. CHECK SHRINKING
		\includegraphics[width=\textwidth]{../../results/report_3D2D_HighLowResCLEAN_meta_exp_075.pdf}
		
		\caption{Multiplot with 2D vs 3D and varying resource densities. The optimal value is denoted by the blue circle.
			The resources used for the low resource scenario (top row) is the minimum amount of resources that allows growth with a $c$ and $\rho$ of 0.01.  As would be expected, since has steeper scaling, 3D allows for growth at smaller resource densities than 2D.
			Low reources in 2D were $ \approx 0.1kg/m^2 $ and $ 0.00035kg/m^2 $ in 3D.
			$ 100kg/m^2 $ was used for high resources because, which although unrealistic, ensures that resources are not a limiting factor in the simulations.
			Hyperallometric scaling is observed in 2D ($\rho = 1.64$ at high resources and 1.99 at low resources (NOTE: This is because 2 is the upper limit of $\rho$ I have simulated here). 
			Scaling in 3D is slightly hypoallometric $rho = 0.94$ and c = 0.01 at high resources. $\rho = 0.8$ and $c = 0.01$ at low resources
			The metabolic scaling exponent = 0.75 in all cases}
		\label{resources2D3D}
	\end{figure}
	
	\begin{figure}[h!]
		
		\includegraphics[width=\textwidth]{../../results/report_2D_exponents.pdf}
		
		\caption{Multiplot with changing metabolic exponents in 2D, ranging from 0.75 to 1, resource density = 0.01, time to maturation is 5 years.  
		It can be seen that higher metabolic responses limit the proportion of mass that can be dedicated to reproduction.  However, in exchange, they allow for much steeper scaling.  A patter which is mirrored in 2D and 3D (fig. \ref{exponents3D})}
		\label{exponents2D}
	\end{figure}
		
	\begin{figure}[h!]
		
		\includegraphics[width=\textwidth]{../../results/report_3D_exponents.pdf}
		
		\caption{Multiplot with changing exponents in 3D, ranging from 0.75 to 1, resource density = 0.01, alpha is 5 years. See fig \ref{exponents2D} for explanation. (NOTE: either 2D or 3D likely belongs in SI) } %either of these likely belong in SI.
		\label{exponents3D}
	\end{figure}
		
	\begin{figure}
		\centering %TODO Units for axis?
		\includegraphics[width=\textwidth]{../../results/pretty_curve}
		\caption{Example of the growth curve and cumulative reproduction expected from a traditional OGM model. Maturation occurs at 1000 days, after which growth is less steep until reaching asymptotic mass.  }
		\label{OGM_Curve}
	\end{figure}
	
	\begin{figure}
		
		\includegraphics[width=\textwidth]{../../results/report_growth_curve.pdf}
		\caption{The growth over a fish which consumes in 2D.  Maturation occurs at 5 years (1825 days).  The fish was allowed to shrink by 5\% at the onset of reproduction.}
		\label{growth_curve}
	\end{figure}



	\nolinenumbers
	
	%%%%%%%%%% Discussion %%%%%%%%%%
\section{Discussion}
	\linenumbers
	% Overall summary of results and their meaning.
	
	
	
	% Issues with my model
		% the mass problem
			% Shrinking and how it addresses the growth speed issue in this model to a degree in my opinion.
		% gsi is likley underestimated in some cases at the very least 0.01 gsi wont be viable for all animals
	Some caveats with the results of the model.  First, growth is simulated as being extremely fast.  As previously stated, asymptotic size was reach by $\sim$10 days.  This is of course not representative of the real world, where individuals generally need several months to years to reach maturity.	
	The rapid growth may be due to several factors.  First is that metabolic cost may be underestimated.  Similar to \citeauthor{West2001} (\citeyear{West2001}), this study used resting metabolic rate to define metabolic costs.  However, this does not take other costs into account such as digestion and locomotion.
	%A rate which is typically used within energetic and growth studies, including early OGMs \parencite{Brett1971,Gillooly2001}  % ref for RMR in growth Tsoukali?  Burger? 
	This was addressed in traditional OGMs by \citeauthor{Hou2008} (\citeyear{Hou2008}).  However, due to the use of asymptotic mass in the parametrisation of this change, the same changes could not be used in this model.  Resting metabolic rate and active metabolic rates do not scale in the same way with mass \parencite{Gillooly2001, Weibel2004}.  The additional cost of active metabolic rate would cause a steeper scaling within the metabolic cost term, leading to more gradual growth (Fig. \ref{scaling_plot}). 
	As such, inclusion of active metabolic rates, while challenging to measure directly and implement, is needed.
	Additionally, there may be behavioural or physiological factors that would lead to an altered metabolic rate.  
	% TODO V presenting this as an issue but it is an area for improvement needs rephrasing and/or shifting
	In this regard, temperature plays a critical role.  It is well documented that a change in temperature will change many biological rates  (see \citeauthor{Peters1983} (\citeyear{Peters1983}), \citeauthor{Gillooly2001} (\citeyear{Gillooly2001}) etc.).  
	It has been shown that growth is dependent on temperature within fish. % sources
	For example, a 1.5\textdegree{}C increase in sea temperatures could result in a 15\% decrease in fish lengths \parencite{VanRijn2017}. % try to phrase better or bolster since it is lifted from diego
	The functional response data used in this study is standardised around 15\textdegree{}C \parencite{Pawar2012}.
	Meanwhile the metabolic cost is for an unspecified temperature \parencite{Peters1983}.  Using rates where the temperature effect is taken into account is crucial for model accuracy given the variation in rates that occurs over different temperatures.  Work such as \citeauthor{Barneche2014} (\citeyear{Barneche2014}) has investigated this effect, however the estimate for metabolic rate is several orders of magnitude lower than what was reported by \citeauthor{Peters1983} (\citeyear{Peters1983}), which is the rate used in this study.  Thus further investigation is required.  % should include a plot in the SI for this.
	
	In fish metabolic rate has also been shown to drop under starvation \parencite{Cook2000}.  In homeotherms feeding restriction has also been shown to also lower body temperature, since metabolic rate and core temperature are closely related in homeotherms \parencite{Ballor1991, Blanc2003,}.  
	
	Another possible point of error is the estimates for supply.  The parameters used from \citeauthor{Pawar2012} (\citeyear{Pawar2012}) are for a spectrum of animals from mammals to insects.  It is possible by reanalysing the data for only marine species, or more specifically only within taxon or species, predictions of supply could be improved \parencite{Marshall2019}.  % check dim paper for addressing rates in different environments
	
	A factor that is not taken into account in this model is that resources are not constant over time.  This can be implemented within the model by varying resource density over time.  The functional response will respond accordingly giving intake which varies through time.  One concern with implementing such a response is fluctuations are likely not experienced by all organisms in the same way.  For a fish with a small range a local fluctuation can be measured and described relatively simply.  However, for a fish with a very large range there is the possiblity leaving resource poor areas in search of richer waters.
	
	
	Despite these caveats, the patterns which arise from the model are promising.  The qualitative patterns seen in this study should not change even with slower growth.  This is because while the exact results are not representative of reality, the relationship between the values is.  For example, the manner in which supply and metabolic rate interact for different values of the metabolic scaling exponent.  The fact that a higher metabolic scaling exponent will cause intake and metabolic cost to intersect sooner does not change regardless of the absolute values.
	
 
	% However such field data is sparse. <- justify saying this some how. need a review or something.
	
	
	% What else can be done to address these issues
		% What is the minimum amount of GSI needed. reproductive structure will vary a bit between species and there is likely a minimum that a species must devote to these structures for them to be viable.
	
	
	% Emphasis that this is applicable to more than just fish - What needs to be done to apply the model to other animal groups.
	In conclusion, the model presented in this study is a promising base which can be expanded upon in a way that was not possible with previous OGMs allowing for much more controlled and detailed explanations of the factors controlling growth.  In contrast to all previous work, which assume optimal supply, the concept of varying supply is addressed using functional responses.  Additionally, qualitative evidence is provided supporting hyperallometric scaling in fish using energy budget as the basis.  The model can easily be applied to any animal taxon, not just fish, with some simple changes.  Additionally, there are clear directions to be explored to improve the model's accuracy.
	
	
	
	
	
	\begin{figure}
		\centering
		\includegraphics[width=\linewidth]{../../results/scalingplot.pdf}
		\caption{a) shows how maintenance cost outscales supply in a traditional OGM.  Growth only stops when maintenance reaches the supply line  b) shows scaling for supply and maintenance as equal. Since scaling is equal growth will never stop until the new cost of reproduction is introduced some time during development.  Blue line is supply, orange line is maintenance, green dotted line is reproduction}
		\label{scaling_plot}
	\end{figure}
%	\subsection{planned outline}
%	Overall summary of results and their meaning.
%	\\- I show that hyper allometry is possible within fish.  
%	\\- higher metabolic cost are leading to hyper allometry, even though I am seeing this for what the data considers too high (i.e. metabolic scaling expontent = 1) \parencite{Marshall2019} % see my google drive doc
%	\\- scaling in 3D will outscale metabolic cost in all cases Table \ref{parameters} because it is greater than 1, need reproductive cost to stop growth.
%	Shrinking and how it addresses the growth speed issue in this model to a degree in my opinion.
%	\\- Is like seeing the stage of growth we expect to see in a typical growth curve in reverse.  Where instead of growing into the asymptotic mass, it shrinks to it.  The energy balance remains the same.
%	\\
%	Issues with my model
%	- mass is too high, addressing a more detailed reproductive mass/metabolic rate may address this or reanalysing samraats data to be more taxon specific may do it.
%	\\
%	What else can be done.
%	\\- more fish specific intake using \cite{Pawar2012} methods as argued in \cite{Marshall2019}
%	\\- trying to model specific species with know growth curves to see how well our results map to them even if the curve isn't great it could be that reproductive output and final mass map well.
%	\\- more detailed metabolic cost since here is basically just RMR with higher exponents.  It may be that the constant shift too and this would lead to better predictions
%	\\- Incorporate oscillating resource density.  was attempted here but due to computational problems could not be implemented.  There is a lack of oscillation data this though and the goal would need to be far more specific since it will depend on scale and range of the organism.
%	\\- This same framework could be applied to other taxa with just a little adjustment (metabolic cost and reproductive equation?)
%	\\- There may be behavioural changes that occur at under feeding restriction that slows metabolism, similar to hibernation etc.
%	\\- temperature could be added here to adjust the metabolic and functional response rates %samraats data is standardised to 15C doesn't seem to be noted in peter's, hemingsen may show it
%	- develop the model to include storage like DEB or put this into a DEB framework.
%	\\
%	What is the minimum amount of GSI needed. reproductive structure will vary a bit between species and there is likely a minimum that a species must devote to these structures for them to be viable. %\cite{Gunderson1997} shows that mortality increases with higher GSI and larger inds have lower GSI does show low values for GSI (0.02) 
%	%GSI can reach 0.75 in invertebrates \cite{Parker2018}
%	\\
%	Emphasis that this is applicable to more than just fish - What needs to be done to apply the model to other animal groups. % need to open up gsi \cite{Parker2018} 0.75 inverts
%	\\
%	This is all energetically based, allowing us to ask questions of why or how there are exceptions to this in the case that organsims should be able to grow larger energetically speaking but may be limited genetically or via some other mechanism
	
	\nolinenumbers
%	\subsection{notes}
%	there are some integration issues that could allow for the appearance of more possible values
%	The shrinking helps to show what could be possible given there is a more gentle growth curve since the main difference is that the time when growth become more limited is just shown as shrinking vs slower growth.


	%%%%%%%%%% Conclusion %%%%%%%%%%
\section{Conclusion}
	\linenumbers
	
	
	
	\nolinenumbers
	
	%%%%%%%%%% Data Declaration %%%%%%%%%%
	\textbf{Code and Data Availability}

	Code is available at: \url{https://github.com/Don-Burns/Masters_Project}
	
	
	%%%%%%%%%% Bibliography %%%%%%%%%%
	\newpage
	\linenumbers
	
	%nicer indentation
	\addcontentsline{toc}{section}{\protect\numberline{}References}	
	
	\printbibliography
	
	% prob more official method of doing the above
%	\printbibliography[heading=bibintoc, title={References}]
	\nolinenumbers
	
	
	%%%%%%%%%% SI %%%%%%%%%%
	\newpage
	%to have table and figures numbered from 1 with prefix
\newcommand{\beginsupplement}{%
	\addcontentsline{toc}{section}{\protect\numberline{}Supplementary Information}
	\setcounter{table}{0}
	\renewcommand{\thetable}{S\arabic{table}}%
	\setcounter{figure}{0}
	\renewcommand{\thefigure}{S\arabic{figure}}%
}

\section*{Supplementary Information}
\beginsupplement

\subsection*{notes}
need section on value conversions and derivations

move any unreferenced sensitivity analyses here.


\begin{figure}
	\centering
	\includegraphics[width=\textwidth]{../../results/Sens_Maturation_Time_Long}
	\caption{Long time scale maturation time}
	\label{fig:sensmaturationtimelong}
\end{figure}

\begin{figure}
	\centering
	\includegraphics[width=\textwidth]{../../results/Sens_Maturation_Time_short}
	\caption{Short time scale maturation time}
	\label{fig:sensmaturationtimeshort}
\end{figure}

\begin{figure}
	\includegraphics[width=\textwidth]{../../results/Sens_Low_Resources_Maturation_Time_short}
	\caption{Low Resources Maturation}
	\label{fig:senslowresourcesmaturationtimeshort}
\end{figure}

\begin{figure}
	\centering
	\includegraphics[width=\textwidth]{../../results/Sens_Metabolic_Exponent}
	\caption{High Resources Metabolic Exponent}
	\label{fig:sensmetabolicexponent}
\end{figure}

\begin{figure}
	\centering
	\includegraphics[width=\textwidth]{../../results/Sens_Low_Resources_Metabolic_Exponent}
	\caption{Low Resource Metabolic Exponent}
	\label{fig:senslowresourcesmetabolicexponent}
\end{figure}

\begin{figure}
	\centering
	\includegraphics[width=\textwidth]{../../results/Sens_Resource_Density}
	\caption{Large scale Resource Resource Density}
	\label{fig:sensresourcedensity}
\end{figure}

\begin{figure}
	\centering
	\includegraphics[width=\textwidth]{../../results/Sens_Low_Resources_Resource_Density}
	\caption{Fine scale resource Density}
	\label{fig:senslowresourcesresourcedensity}
\end{figure}

\begin{figure}
	\centering
	\includegraphics[width=\textwidth]{../../results/Sens_Shrink}
	\caption{High Resources Shrinking}
	\label{fig:sensshrink}
\end{figure}

\begin{figure}
	\centering
	\includegraphics[width=\textwidth]{../../results/Sens_Low_Resources_Shrink}
	\caption{Low resources shrinking, missing values around 0 are because no viable solutions could be found.}
	\label{fig:senslowresourcesshrink}
\end{figure}




	


	
	
\end{document}} % to include word count 

%%%%%%%%%% Formatting %%%%%%%%%%
\usepackage[margin=2cm]{geometry} % margins of 2cm
\linespread{1.5} %1.5 spacing
%\renewcommand{\familydefault}{\sfdefault} % set font to arial clone (helvet)

\usepackage[compact]{titlesec} % reduce spacing bewteen section titles

%%%%%%%%%% Bibliography %%%%%%%%%%
\addbibresource{../Masters_Thesis.bib}

%commands 

%%%%%%%%%% Document %%%%%%%%%%
\begin{document}
	%%%%%%%%%% Title Page %%%%%%%%%%
	%\newcommand{\crest}{\includegraphics[width = 4cm, keepaspectratio]{../images/IC_Crest.eps}} % Imperial crest
%%formating
\begin{titlepage} % Suppresses headers and footers on the title page
	
	\centering % Centre everything on the title page
	
	\scshape % Use small caps for all text on the title page
	
%	\vspace*{\baselineskip} % White space at the top of the page
	
	%------------------------------------------------
	%	Title
	%------------------------------------------------
	
	\rule{\textwidth}{1.6pt}\vspace*{-\baselineskip}\vspace*{2pt} % Thick horizontal rule
	\rule{\textwidth}{0.4pt} % Thin horizontal rule
	
	\vspace{0.75\baselineskip} % Whitespace above the title
	
	{\LARGE Using Realistic Resource Supply to Show Hyperallometric Reproductive Output in Fish\\} % Title
	
	\vspace{0.75\baselineskip} % Whitespace below the title
	
	\rule{\textwidth}{0.4pt}\vspace*{-\baselineskip}\vspace{3.2pt} % Thin horizontal rule
	\rule{\textwidth}{1.6pt} % Thick horizontal rule
	
	\vspace{1\baselineskip} % Whitespace after the title block
	
	%------------------------------------------------
	%	Subtitle
	%------------------------------------------------
	
	%SUBTITLE? % Subtitle or further description
%	Student:
	
	
	\vspace{0.5\baselineskip} % Whitespace before 
	
	{\scshape\Large D\'onal Burns  \\} % my name
	
	\vspace{0.5\baselineskip} % Whitespace below 
	
	\textit{CID: 01749638 \\ Imperial College London \\ Email: donal.burns@imperial.ac.uk} % affiliation and email
	
	\vspace*{2\baselineskip} % Whitespace under the subtitle
	
	
	
%	Supervisor:
%	
%	
%	\vspace{0.5\baselineskip} % Whitespace before 
%	
%	{\scshape\Large Samraat Pawar \\} % supervisor name
%	
%	\vspace{0.5\baselineskip} % Whitespace below 
%	
%	\textit{Imperial College London \\ Email: s.pawar@imperial.ac.uk} % affiliation and email
%	
%	\vspace{3cm} % Whitespace between 
	

	
	%% crest
	
	\includegraphics[width = 4cm, keepaspectratio]{../images/IC_crest.pdf}
	
	\vspace{0.3\baselineskip} % Whitespace under the Uni logo
	
	Submitted: August 27$^{th}$ 2020 % Publication Date
	
	

		%%Submission clause
	\vspace{0.5\baselineskip}
	
	A thesis submitted in partial fulfilment of the requirements for the degree of
	%Computational Methods in Ecology and Evolution 
	Master of Science at Imperial College London
	\vspace{0.5\baselineskip}
	
	Formatted in the journal style of Functional Ecology	
	\vspace{0.5\baselineskip}
	
	Submitted for the MSc in Computational Methods in Ecology and Evolution
	
\end{titlepage}

	

	%%%%%%%%%% Declaration %%%%%%%%%%
	\section*{Declaration}
	I declare this project as my own work.  The model presented here was developed in conjunction with my supervisor, Dr. Samraat Pawar, and Ph.D. students Tom Clegg and Olivia Morris.  I was responsible for any simulations and data presentation.\newline
	%% word count
	\textbf{Word Count: \wordcount}

	\newpage
	
	%%%%%%%%%% Abstract %%%%%%%%%%
	\section*{Abstract}
	\linenumbers
%	Size is essential to reproductive output. By extension understanding growth, determines size allows understanding of reproductive output. 
	 
	With recent results showing that reproduction in fish scales hyperallometrically there is a need to update growth OGMs to reflect this fact.  Current OGMs assume optimal intake, an assumption which is not always reflected in the field.  This study develops an energy intake focused approach to explaining growth, an area which has not been covered within current literature, and shows that hyperallometric scaling of reproductive output arises when allowing for variable reproductive scaling and maximising for fitness.  The model is applicable to not only fish, but any animals taxon with some simple parameter adjustments.  I offer direction for improvements and areas to be developed in order to allow the model to be applicable to any temperature range.
	\vspace*{0.5 cm}
	\newline
%	TODO pieces to include:
%	
%	show the effects of the exponent and why it is important that we use the correct one moving forward.
%	give direction as to what needs to be done going forward.  what areas need work
	%TODO KEywords
	\textbf{Keywords:}\\
	allometry; functional response; metabolic theory; growth; intake; life history; metabolism; reproduction; reproductive output; supply

	
	
	\nolinenumbers
	%%%%%%%%%% Acknowledgements %%%%%%%%%%
	%\thispagestyle{empty}

\mbox{}\newline\vspace{10mm} \mbox{}\LARGE
%
{\bf Acknowledgements} \normalsize \vspace{5mm}

%I would like to thank my supervisor.....
%thanks to luke who gave a great base to start from on the project





	
	%%%%%%%%%% Table of Contents %%%%%%%%%%
	\tableofcontents
	\newpage

	%%%%%%%%%% Introduction %%%%%%%%%%
\section{Introduction}
	\linenumbers
	% putting OG first
	Growth is a gateway to understanding how many biological rates change throughout ontogeny.  This is because so many biological rates scale with mass \parencite{Kleiber1932}.  To date many models have been developed to predict and describe the growth of an organism throughout its lifetime.  The three main approaches used are the von Bertalanffy model, the dynamic energy budget (DEB) model and the ontogenetic growth model (OGM), which will be the focus of this study \parencite{Putter1918, vonBertalanffy1938, Kooijman1986, West2001}.  All of these are energetic based models with varying assumptions, key among which is the scaling of resource supply and metabolic rate with mass.   In OGMs supply is thought of as being optimal at all times which leads to the assumption that intake scales to the power of 0.75.  Indeed while under optimal conditions this may be true, however, it neglects that this situation is thought to rarely occur in the field \parencite{Pawar2012}.
%	Additionally, growth model approaches have made various assumptions about the scaling of metabolic rates.  OGMs assume that scaling is isometric with mass, i.e. scaling exponent is 1.  This is contrary to what has been thought 
	This study takes the novel approach of using a mass specific functional response and assimilation efficiency to describe how intake changes both throughout ontogeny and varying levels of resource availability. While the methodology used here can be used for any animal taxon, in order to investigate growth, a taxon much be chosen to parametrise metabolic rate around.  This study will focus on fish.
		
% TODO bring in OGM growth curve and use to illustate the basic OGM as per samraats comments
	% Why is growth important
	Understanding how organisms grow and what factors play a role in determining growth is not only insightful from the perspective of understanding the world around us, but can also be used to better manage the many fisheries and marine protected areas around the world.  An objective which is becoming increasingly important as the oceans' fish stocks continue to be depleted by over fishing. 
	%VV rephrase this part to hint more towards size than it currently does VV
	This is compounded by global warming which threatens to alter the structure of the ecosystems even if they are not fished and left in their ``natural" state. 
	
	Larger fish produce more offspring than smaller ones.  % this is known in general isometrically
	In fact, it has been shown that larger fish produce far more offspring than the equivalent mass composed of smaller fish,  that is a 2kg fish will produce more offspring than two 1kg fish \parencite{Barneche2018}.
	Larger fish also use energy more efficiently than multiple smaller ones per unit mass.  This is due to them having a lower mass specific metabolic rate \parencite{Peters1983, Kleiber1932, Brown2004}.  
	%As a result a higher proportion of larger fish may actually reduce stress upon the ecosystem from an energetic perspective which could be useful when trying to manage fish stocks. %need citations for this also questionably neccessary as a point and has little backup? may even not be the case at all due to other pressures introduced by larger individuals in a system
	% larger mothers produce larger offspring which may better survive \cite{Barneche2018}
	Additionally, larger mothers produce larger offspring, which are then more likely to survive to adulthood and reproduce \parencite{Hixon2014, Marshall2006}. % double check this statement from barneche and get more backing
	%
	
	It is already known that metabolic rate is dependant on temperature which in turn will lower fish sizes \parencite{Gillooly2001, Brown2004} and with increasing global temperatures, understanding in greater detail how increased metabolic rates may affect growth is useful in population management.
	%
	
%	Growth is in essence a balancing act between how much energy an organism can acquire and the amount of energy required for maintenance, 
%	that is movement, digestion et cetera.  
	
	%power laws -- needs citation or could be counted as common knowledge? \cite{Huxley1932}
	Key to understanding biological rates and their relationship with mass is the concept of power laws. Many biological traits can be described as scaling to the power of some other biological trait.  That is some rate, $Y$, can be expressed for any mass by $Y = Y_0 m^\beta$, where $m$ is mass \parencite{Kleiber1932}.  Power laws can be broadly categorised based on the value of their exponent, $\beta$.  
	Where the exponent equals one the relationship is described as isometric or linear, that is the two traits increase at the same rate.  
	Where the exponent does not equal one, the relationship is said to be allometric.  That is the trait does not increase at the same rate as the trait being compared against.  In cases where the exponent is greater than one, the trait is said to scale super-linearly or hyperallometrically and to scale sub-linearly or hypoallometrically for cases with an exponent of less than one.
	For the purposes of describing growth these relationships are indispensable.   
	
	
%	Traditionally ontogenetic growth models have relied on knowing how large an organism is expected to grow.
%	%	von Bertalanffy / puetter model
%	The von Bertalanffy growth equation relies on knowing the longest a fish can be and the length of the fish at the beginning of measurements \parencite{vonBertalanffy1938, Putter1918}.  From here, for a known growth rate, the length of the fish after a set amount of time has passed can be predicted.
%	
%	%OGMs
%	Another approach has been to try explain growth from an energetic perspective.  That is an individuals growth is determined by the amount of energy available to it.  Once all costs, such as maintenance, have been paid all remaining energy can be invested towards growth.  There are two frameworks which have been built around this concept; dynamic energy budget (DEB) theory % TODO get ref for DEB original/important sources %add note to see discussion for more on DEB?
%	and ontogenetic growth models (OGMs), which will be focused on here.

	%	introduce \cite{West2001}	
	One of the best known examples of an OGM is the model developed by \cite{West2001}.  This model is parametrised around the average energy content of animal tissue and asymptotic mass.  Asymptotic mass being the mass at which growth has essentially stopped due to metabolic cost and energy intake equalling each other (Fig. \ref{scaling_plot}a). The model hinges on the scaling with mass between energy intake (m$^{0.75}$, i.e. allometric sub-linear scaling) and maintenance cost (m$^1$, i.e. isometric linear scaling).  In other words, as mass increases maintenance costs will slowly overtake the intake rate and halt growth.  	
	
	% talk about determinant and indeterminant growth/
	% move to including reproduction with \cite{Charnov2001} mashed with \cite{Hou2008} imporvements briefly
	The framework used by \citeauthor{West2001} (\citeyear{West2001}) was latter developed by \citeauthor{Charnov2001} (\citeyear{Charnov2001}) to take the cost of reproduction into account and allowing the estimation of lifetime production of offspring.  \citeauthor{Hou2008} (\citeyear{Hou2008}) latter developed \citeauthor{West2001}'s model further by expanding maintenance cost to include the cost of feeding and digestion (specific dynamic action), synthesis of new tissue and activity.
	
%	% move to including reproduction with \cite{Charnov2001}
%	The framework used by \citeauthor{West2001} (\citeyear{West2001}) was latter developed by \citeauthor{Charnov2001} (\citeyear{Charnov2001}) to take the cost of reproduction into account and allowing the estimation of lifetime production of offspring.  This was done by adding an additional cost which dedicated some portion of mass to reproductive output, essentially gonadosomatic index.	%like 20% max? in fish very high in  invertabrates \cite{Parker2018}
%	% improved by hou et al \cite{Hou2008} talk about including SDA and energy improvements
%	\citeauthor{West2001}'s (\citeyear{West2001}) model was later expanded upon by \citeauthor{Hou2008} \citeyear{Hou2008} to elaborate on the maintenance cost.  \citeauthor{Hou2008} (\citeyear{Hou2008}) takes both the difference between the cost of maintaining tissue and creating it, a factor that is neglected in \citeauthor{West2001} (\citeyear{West2001}), and the cost of digestion into account.
	
	% begin caveats
	% tautology still present 

	% Discuss allometry and isometry here to highlight what scaling super or sub linearly means
	In the above OGMs intake is assumed to scale sub-linearly to the power of 0.75.  This is due to the assumption that the individuals are consuming at an optimal rate at all times and therefore the only limitation is their ability to make use of that energy.  In this case, intake should theoretically scale to the power of 0.75 (see \cite{West1997}).  However, this is not always the case in the field.  It has been shown that, for non-optimal consumption, steeper scaling can occur \parencite{Pawar2012, Peters1983}. Additionally, OGMs, like many growth and metabolic models, typically use basal or resting metabolic rate to calculate metabolic cost.  This is the minimal metabolic rate of an organism and is typically thought of as the rate of the organism when at relaxed and at rest.  However, it has been shown, once factors such as movement are taken into account, that the scaling becomes steeper \parencite{Weibel2004}.
	% still no accounting for imperfect intake though is somewhat covered by \cite{Hou2011}
	% bring up the caveats/assumptions of the model i.e. isometric metabolism and repro, 0.75 scaling of intake. optimal intake at all times etc. 
	% tautology of OGMs and \cite{Hou2011} close but still issues
	The issue of non-optimal feeding is addressed somewhat by \citeauthor{Hou2011} (\citeyear{Hou2011}).  However, this growth was only investigated as, essentially, a proportion of optimal consumption and does not address a potential change in scaling of intake rate.
	Another limitation of the models used in previous OGMs is the dependence on asymptotic mass.  The models are entirely dependant on the value of optimal intake and asymptotic mass.  All other values, such as metabolic cost, are then derived in relation to these.  However, organisms are not born with some inherent restriction to the size they can attain, at least not energetically.  If there is surplus energy for a given mass the organism should be able to grow.  Relying on asymptotic mass to define the upper bound of attainable mass does not allow for investigation of the mechanisms that underpin asymptotic size in reality. 
	
	% talk about how we assume that evolution settles upon values which optimise for reproduction -- can do this when talking about charnov?
	
	% Justify and appeal my methods	
	With two of the key assumptions of current OGMs, that reproduction and metabolism scale isometrically, not holding in the field \parencite{Barneche2018, Pawar2012, Peters1983} there is a need to take an unexplored approach to modelling fish growth, in particular choosing to focus on developing how intake is described so as to better reflect the real world.  To do this an obvious starting point is to model intake as a functional response \parencite{Holling1959} so as to better reflect real world intake rates in terms of consumed biomass over time.  Non-optimal supply is a currently unexplored area within growth modelling.  This is likely due to the difficulty of directly measuring intake, especially in the field, Perhaps as a result comparatively less is known about consumption.  This leads to the need for proxy values to estimate it, for example nutrient flux \parencite{Schiettekatte2020}, or drawing broad relationships to approximate consumption, as this study will do.
	Changing the manner in which intake is defined also requires changing metabolic cost, because the two are dependent upon each other in current OGMs.  This can be done by defining metabolic rate as a value dependent on current mass rather than asymptotic mass as has been done in OGMs up until this point.  This thought process is more mechanistic since an organism has no concept of ``How large should I grow?", but rather will acquire as much resources as it capable of at its current life stage and size.  Taking this more bottom-up mechanistic approach also allows exploration of what controls how large fish grow to, since as previously mentioned, from an energetic standpoint, an organism can grow indefinitely as long as there is surplus energy available after costs have been paid.  Of course, there are also mechanical and genetic limitations upon organism size, but once size is constrained to what is known to exist, this should not be an issue.  
	Under the assumption that fish have evolved to maximise reproductive output, using the above framework, simulations can be carried out to demonstrate what conditions need to be met in order to achieve hyperallometric scaling of reproduction from an energetic perspective.  This study will answer 1) Under what metabolic scaling exponent can hyperallometric reproduction arise, 2) 
	% describe basic OGM and how i change the intake term to reflect supply
	
%	\cite{Sibly2020} % for metabolic scaling
%	\cite{Marshall2019} non-isometric scaling of repro#

% questions
% what effect does exponent actually have in relation to growth
% what coditions need to be fuflilled to possibly acheive hyperallometry
% 
	\nolinenumbers
	
	%%%%%%%%%% Methods %%%%%%%%%%
\section{Methods}
	\linenumbers
	
	\subsection{Altering OGMs to account for supply}
	In order to address the issue of supply in the context of an OGM, which can be generically described as $dm/dt = gain - loss$, some changes need to be made to the model's terms.  The first is to remove the assumption of asymptotic mass and the reliance of cost upon it.  Within a tradional OGM the gain term ($ a $) and asymptotic mass are used to define the metabolic cost ($ b $).  However, since the assumption of perfect intake is going to be broken, because of the variable supply, this relationship no longer holds.  As such, both intake and metabolic cost need to be redefined.  Additionally, in light of recent work showing that reproduction scales allometrically and not isometrically, the reproductive cost must also be modified from the form used by \citeauthor{Charnov2001} (\citeyear{Charnov2001}) \parencite{Marshall2019, Barneche2018} %TODO check marshal ref data origin
	
	
	\subsubsection{Gain}
	To define intake a natural starting place is the functional response \parencite{Holling1959}.  Functional responses  are used to define how much an organism consumes for a given resource density and is described by the following equation:	
	\begin{equation}
		\label{functional_repsonse}
		f(\cdot) = \frac{a X_r}{1 + a h X_r}
	\end{equation}
	where, $ f(\cdot) $ is the functional response, $ a $ is the search rate, $ h $ is handling time and $ X_r $ is resource density.  
	For a fixed mass and increasing resource density Eq \ref{functional_repsonse} produces a sigmoidal shape with intake eventually reaching an asymptote after some saturating amount of resources is reached.  At lower resource densities, the intake rate is primarily defined by the search rate with higher search rates yielding higher intake rates.  Conversely, at high resource densities, intake rate is approximately equal to the inverse of the handling time ($ h^{-1} $), where lower handling times yield higher intake rates.  
	
	An organism's functional response will not remain constant throughout its life history.  Search rate and handling time are affected by both the organism's mass and how it interacts with its environment \parencite{Pawar2012}.  
%	Within this model mass will be known for all time points since that is one of the quantities being predicted.  
	Interactions can be broken into 3D and 2D, that is whether the organism consumes from a 2D ``surface" e.g. a cow grazing or a 3D ``volume" e.g. a pelagic consumer which consumes prey from within the water column.  As such, both handling time and search rate can be defined as Eq. \ref{handling_time} and Eq. \ref{search_rate} respectively.
	\begin{equation}
		\label{search_rate}
		a(m) = a_0 m_t^\gamma
	\end{equation}
	
	\begin{equation}
		\label{handling_time}
		h(m) = t_{h,0} m_t^\beta
	\end{equation}
	A functional response alone is not enough to fully define intake.  This is because processing of consumed resources is not one hundred percent efficient which leads to inevitable loss consumed energy.  As a result, to achieve the final gain term, a dimensionless efficiency term $\epsilon$ is applied.  In poikilotherms assimilation efficiency is roughly 70\% \parencite{Peters1983}
	
	\subsubsection{Loss}
	Metabolic cost  has previously been dependant upon the gain term within traditional OGMs (see \cite{West2001, Hou2008}).  However, for non-maximal intake the relationship will no longer hold true.  As a result, this model has taken previously measured values to be used as metabolic cost (see Eq. \ref{metabolic_cost} taken from \cite{Peters1983} % add ref to hemmingsen here since that is the peters source
	and Table \ref{parameters} for further details)
	\begin{equation}
		\label{metabolic_cost}
		B_m = 0.14 m_t^\mu
	\end{equation}
	Next to take allometric scaling of reproduction into account, the reproductive cost term from \citeauthor{Charnov2001} (\citeyear{Charnov2001}) is changed from $cm^1$ which assumes isometric scaling to $cm^\rho$.  $c$ can be interpreted as the proportion of mass dedicated to reproduction, i.e. the gonadosomatic index (GSI) of the fish \parencite{Charnov2001}.  Just as in \citeauthor{Charnov2001} (\citeyear{Charnov2001}) reproductive cost is only taken into account once maturity is reached.  Meaning that until a length of time ($\alpha$) has passed, reproductive cost is considered to be zero.
	
	\subsubsection{Full Growth Equation}
	When all newly defined terms are taken together they produce Eq. \ref{dmdt_juvenile}  and \ref{dmdt_mature} which describe the change in mass for a fish before (Eq. \ref{dmdt_juvenile}) and after maturation (Eq. \ref{dmdt_mature}).
	\begin{align}
		\label{dmdt_juvenile}
		\frac{dm}{dt} &= \epsilon f(\cdot) - B_m & m_t < m_\alpha \\
		\label{dmdt_mature}
		\frac{dm}{dt} &= \epsilon f(\cdot) - B_m - cm_t^\rho & m_t \geq m_\alpha
	\end{align}

	
	
	\subsection{Reproductive Output}
	At any time ($ t $) a reproducing organism devotes some amount of energy to reproduction.  This is the product between the amount of mass dedicated to reproduction ($ cm^\rho $) and a declining efficiency term ($ h_t $) which begins at maturity ($ \alpha $) and represents reproductive senescence \parencite{Benoit2018, Vrtilek2018, Stearns2000}.  In addition to amount of reproduction, the offspring are also subject to mortality ($ l_t $).  By combining the two, lifetime reproductive output can be estimated and is described by the ``characteristic equation" (Eq. \ref{characteristic_equation}) which represents reproductive output in a non-growing population \parencite{Tsoukali2016, roff1993, Roff2001, stearns1992evolution, Arendt2011, Roff1986, Roff1984}
	\begin{equation}
		\label{characteristic_equation}
		R_0 = \int c m_t^\rho h_t l_t 
	\end{equation}
	Mortality is experienced differently by juvenile ($ t < \alpha $) and reproducing individuals ($ t \leq \alpha $). % TODO <-dont like this sentence
	Mortality of offspring prior to maturity is described as a survival rate $ l_t = e^{-Z(t)} $ which is an exponentially decreasing function bounded between zero and one.  It controls how many offspring make it to maturity.  After maturity, survival is again described as an exponential function which takes time to maturity into account, $ l_t = e^{-Z(t-\alpha)} $.  
	Reproductive senescence can be also be estimated as an exponential function which begins after maturity and declines over time  ($ e^{-k(t-\alpha)} $), where $ k $ is the senescence term.  When all values are inserted into the characteristic equation, it results in the equation used by \citeauthor{Charnov2001} (\citeyear{Charnov2001}) with the inclusion of reproductive senescence (Eq. \ref{reproductive_output}).
	\begin{equation}
		\label{reproductive_output}
		R_0 = c\int_0^\alpha e^{-Z_t} dt  \int_\alpha^\infty m_t^\rho e^{-(\kappa + Z)(t - \alpha)}dt\\
	\end{equation} 
	In Eq. \ref{reproductive_output} $ Z $ represents instantaneous mortality.  This rate has been shown to be related to time of maturation in many taxon groups, and within it follows the relationship $ \alpha \cdot Z \approx  2$.  This can then be rearranged to estimate instantaneous mortality, $ Z \approx 2/\alpha  $
	
	\subsection{Maximising Reproduction}
	It is assumed that evolution will converge on metabolic values which maximise fitness. 
	Fitness being defined as how much an individual is able to contribute to the gene pool \parencite{Speakman2008, Stearns2000}.  % Would like to remove this and link more smoothly
	To this end, lifetime reproductive output is often used as a measure of fitness \parencite{Charnov1991, Audzijonyte2018, Speakman2008, Stearns2000, Charnov2001, Tsoukali2016, Brown1993, Charnov2007}.  Therefore, by maximising for reproductive output, it should become clear what parameters will yield the highest fitness.  These parameters will then show whether, within a theoretical framework, hyperallometric scaling arises.
	
	To find all optimal values for reproduction would require Eq. \ref{reproductive_output} to analytically solved.  However, since no such solution is possible, I simulated the problem numerically to obtain a result.  This was done by simulating across values of $ c $ and $ \rho $, the parameters of interest between growth (Eq. \ref{dmdt_juvenile} and \ref{dmdt_mature}) and reproductive output (Eq. \ref{reproductive_output}).  $ c $ was bound between 0 and 0.4, which encapsulates the values measured within fish \parencite{Benoit2018, Roff1983, Fontoura2009}.  Though it has been shown to reach as much as 0.7 in some species \parencite{Parker2018}  To search for any hyperallometry within reproduction, $ \rho $ was bound between 0 and 2.  
	The simulation was then run at 0.01 value intervals in both $c$ and $\rho$ over a lifespan of one million days.  The results of each simulation were recorded and any non-viable results were discarded.  A result was considered non-viable if fish had ``shrunk" more than 5\% in order to accommodate reproductive costs.  Shrinking occurs in the model because  the combined loss of energy to metabolism and reproduction is too much for the simulated values at the mass achieved by maturation so the individual experiences a deficit of energy, which is paid by loss in mass until equilibrium is achieved. % \cite{VandenBerghe1992} for reproductive mass loss (though it is due to behaviour changes)
	Shrinking is not expected at maturity in reality.  Typically, maturity will occur while the organism still has room for growth.  It is the onset of reproduction which is considered to slow or stop growth % this is the case where the metabolic exponent is the same or less than the intake one.
	 (see Fig. \ref{OGM_Curve}).  Shrink can be thought of as starvation in a real organism.  If energetic cost are not met then energy reserves in the body, such as fat and muscle, are broken down for energy.  It has been shown that some fish are capable of losing up to 10\% of their body mass \parencite{VandenBerghe1992}.  However, this was during the breeding season and caused by behavioural changes due to parenting.  Additionally, individuals were shown to rebound back to their``normal" body mass once the breeding season had ended.
	% survival impacts of shrinking? are there sources?
	% \cite{Lambert2000}
	
	\subsection{Impact of Parameters}

	In order to determine the roles of metabolic exponent, maturation time and resource density within the model, sensitivity analyses were performed on each parameter with regard to $c$ and $\rho$.  This was done by simulating the parameters across multiple values and obtaining the optimal value for $c$ and $\rho$ as described above.
	
	
	\begin{centering}
		
		%		TODO check description capitalisation
		
		\begin{table}[h!]
			\caption{Table describing parameters used in the model, along with values, units and sources where applicable.  The units of resource density change depending on the dimension of intake.  $m^D$ represent either $m^2$ in 2D or $m^3$ in 3D} 
			\label{parameters}
			\vspace{2mm}
			\begin{tabular}{c l l l l p{3cm}}
				\hline
				Parameter 	& Description 			& Value 	& Units 	& Range 		& Source \\
				\hline
				$m$			& Mass					& -			& kg day$^{-1}$& -			&		\\
				
				$B_m$		& Metabolic Cost		& $0.14 m^{\mu}$ & kg day$^{-1}$& - 	& \cite{Peters1983}\\
				$\mu$		& Metabolic Exponent	& -			&	-		& 0.75 - 1.0	& - \\
				$\alpha$	& Age of maturity		& -     	& day		& -				& -\\
				$c$			& Reproduction scaling constant & - & kg day$^{-1}$& 0 - 0.5 		& -\\
				$\rho$		& Reproduction scaling exponent	& -	&	-		& 0 - 1.5			& -\\
				$Z$			& Rate of instantaneous mortality& $2/\alpha$	& & & \cite{Charnov2001}\\%double check ref
				$k$			& Reproductive senescence & 0.01\\
				
				$\epsilon$	& Assimilation Efficiency & 0.70 & - & - 		& \cite{Peters1983} \\
				$X_r$ 		& Resource Density		& -		& kg/m$^D$		& 0.11 - 30				& -\\
				$\gamma$	& Search rate scaling exponent & 0.68 in 2D	& - & - & \cite{Pawar2012} \\
				&						& 1.05 in 3D\\
				$a_0$		& Search rate scaling constant & $10^{-3.08}$ in 2D & m$^2$ s$^{-1}$ kg$^{-0.68}$   & - &\cite{Pawar2012}	\\
				&						& $10^{-1.77}$ in 3D& m$^2$ s$^{-1}$ kg$^{-1.05} $\\
				$\beta$		& Handling time scaling exponent& 0.75 & - & - & \cite{Pawar2012}\\
				$t_{h, 0}$	& Handling time scaling constant& $10^{3.95}$ in 2D &kg$^{1-\beta}$ s& -& \cite{Pawar2012}	\\
				&						& $10^{3.04}$ in 3D			&kg$^{1-\beta}$ s\\
				\hline
			\end{tabular}
		\end{table}
	\end{centering}
%	\subsection{notes}

	\nolinenumbers
	%%%%%%%%%% Results %%%%%%%%%%
\section{Results}
	\linenumbers
	
%	\subsection{notes}
%		scenesence not much of a factor
% 10 \% predicted under Roff 1983

	\subsection{Produced Growth Curve and Maturation Time}
	% talking about at meta_exp = 1
	As can be seen in Figure \ref{growth_curve}, growth is very fast within the model.  Asymptotic mass is reached by $\sim$15 days.  This makes interpreting any results regarding maturation time difficult % TODO < rephrase
	because any time after $\sim$15 days produces the same result (see Fig. \ref{fig:senslowresourcesmaturationtimeshortexp1} - \ref{fig:senshighresourcesmaturationtimeshortexp1}).
	However, as can be seen in Fig. \ref{fig:sensmaturationtimeshortexp075}, where maturation is occurring early during the growth phase, hyperallotric scaling emerges.  This pattern shows in particular at lower resource densities in 3D (Fig. \ref{fig:sensverylowresourcesmaturationtimeshortexp075} and \ref{fig:sensverylowresourcesmaturationtimeshortexp1})
	
	
	\subsection{Effect of Parameters}
	\subsubsection{Resource Density}
	The scaling relationship emerges as would be expected from the scaling of the functional response.  At low resource densities the output of the functional response will scale similarly to search rate.  The scaling of which is higher in 3D as resources increase and the response shifts to scaling similarly to the inverse of handling time.  At this point $\rho$ starts to take values which are higher in 2D than 3D, because of the higher normalisation constant in 2D.
	
	\subsubsection{Metabolic Exponent}
	% talking about at meta_exp = 1
	The expected result is for increasing values of metabolic scaling exponent ($\mu$) for $\rho$ to also increase.  This is because the lower values of $\mu$ will result in there being a larger gap between the scaling of intake and maintenance, which allows for steeper scaling in reproduction (see Fig. \ref{scaling_plot} b)) \parencite{Marshall2019}.  This appears to be the case when analysing $\rho$ with respect to the other parameters where $\mu = 1$ or 0.75 (e.g. Fig. \ref{fig:senshighresourcesmaturationtimeshortexp1} and \ref{fig:sensmaturationtimeshortexp075} or Fig. \ref{fig:sensresourcedensityexp1broad} and \ref{fig:sensresourcedensityexp075fine}).  However, when $\mu$ is explicitly simulated over a variety of values, the trend suggests that increasing $\mu$ allows for higher values of $\rho$.  
	
	\subsubsection{Shrinking}
	% talking about at meta_exp = 1
	At saturated resource densities, allowing for greater proportions of shrinking enables $\rho$ to be set to larger values in both dimensions.  This is perhaps not surprising as increasing the shrinking proportion allows for a larger reproductive cost since all individuals will already be at asymptotic mass at time of maturation.
	
	\subsubsection{$c$ Values}
	Estimations of $c$ are low in many cases, especially in 2D.  While this may be low compared to the $\sim$10\% - 35\% expected \parencite{Benoit2018, Fontoura2009, Roff1983} it is not unprecedented for values of 2\% to be observed in some fish \parencite{Gunderson1997}.  It may be necessary for the lower bounds of $c$ to be adjusted based on what is expected or even viable in the organisms being simulated.
	
%	minimum viable resource is dependent upon meta.exp, with 0.75 allowing for less than that of 1 and scaling in between 0.2 for 1 and 0.15 for 0.75. likewise 3D is capable of going lower than 2D due to the higher intake rate.  shrinkage plays no role here.

	\begin{figure}[h!]
				
%		TODO 1.The title for each plot needs to be cleaned up.  2.The units should be m^2 and 3.need letter labels for each.
		\includegraphics[width=\textwidth]{../../results/report_3D2D_HighLowResCLEAN_meta_exp_1.pdf}
		
		\caption{Multiplot with 2D vs 3D and varying resource densities. The optimal value is denoted by the blue circle.
		The resources used for the low resource scenario (top row) is the minimum amount of resources that allows growth with a $c$ and $\rho$ of 0.01.  As would be expected, since has steeper scaling, 3D allows for growth at smaller resource densities than 2D.
		Low reources in 2D were $ \approx 0.1kg/m^2 $ and $ 0.00035kg/m^2 $ in 3D.
		$ 100kg/m^2 $ was used for high resources because, which although unrealistic, ensures that resources are not a limiting factor in the simulations.
		Hyperallometric scaling is observed in 2D ($\rho = 1.64$ at high resources and 1.99 at low resources (NOTE: This is because 2 is the upper limit of $\rho$ I have simulated here). 
		Scaling in 3D is slightly hypoallometric $rho = 0.94$ and c = 0.01 at high resources. $\rho = 0.8$ and $c = 0.01$ at low resources
		The metabolic scaling exponent = 1 in all cases}
		\label{resources2D3D}
	\end{figure}

	\begin{figure}[h!]
		
		%		TODO 1.The title for each plot needs to be cleaned up.  2.The units should be m^2 and 3.need letter labels for each. CHECK SHRINKING
		\includegraphics[width=\textwidth]{../../results/report_3D2D_HighLowResCLEAN_meta_exp_075.pdf}
		
		\caption{Multiplot with 2D vs 3D and varying resource densities. The optimal value is denoted by the blue circle.
			The resources used for the low resource scenario (top row) is the minimum amount of resources that allows growth with a $c$ and $\rho$ of 0.01.  As would be expected, since has steeper scaling, 3D allows for growth at smaller resource densities than 2D.
			Low reources in 2D were $ \approx 0.1kg/m^2 $ and $ 0.00035kg/m^2 $ in 3D.
			$ 100kg/m^2 $ was used for high resources because, which although unrealistic, ensures that resources are not a limiting factor in the simulations.
			Hyperallometric scaling is observed in 2D ($\rho = 1.64$ at high resources and 1.99 at low resources (NOTE: This is because 2 is the upper limit of $\rho$ I have simulated here). 
			Scaling in 3D is slightly hypoallometric $rho = 0.94$ and c = 0.01 at high resources. $\rho = 0.8$ and $c = 0.01$ at low resources
			The metabolic scaling exponent = 0.75 in all cases}
		\label{resources2D3D}
	\end{figure}
	
	\begin{figure}[h!]
		
		\includegraphics[width=\textwidth]{../../results/report_2D_exponents.pdf}
		
		\caption{Multiplot with changing metabolic exponents in 2D, ranging from 0.75 to 1, resource density = 0.01, time to maturation is 5 years.  
		It can be seen that higher metabolic responses limit the proportion of mass that can be dedicated to reproduction.  However, in exchange, they allow for much steeper scaling.  A patter which is mirrored in 2D and 3D (fig. \ref{exponents3D})}
		\label{exponents2D}
	\end{figure}
		
	\begin{figure}[h!]
		
		\includegraphics[width=\textwidth]{../../results/report_3D_exponents.pdf}
		
		\caption{Multiplot with changing exponents in 3D, ranging from 0.75 to 1, resource density = 0.01, alpha is 5 years. See fig \ref{exponents2D} for explanation. (NOTE: either 2D or 3D likely belongs in SI) } %either of these likely belong in SI.
		\label{exponents3D}
	\end{figure}
		
	\begin{figure}
		\centering %TODO Units for axis?
		\includegraphics[width=\textwidth]{../../results/pretty_curve}
		\caption{Example of the growth curve and cumulative reproduction expected from a traditional OGM model. Maturation occurs at 1000 days, after which growth is less steep until reaching asymptotic mass.  }
		\label{OGM_Curve}
	\end{figure}
	
	\begin{figure}
		
		\includegraphics[width=\textwidth]{../../results/report_growth_curve.pdf}
		\caption{The growth over a fish which consumes in 2D.  Maturation occurs at 5 years (1825 days).  The fish was allowed to shrink by 5\% at the onset of reproduction.}
		\label{growth_curve}
	\end{figure}



	\nolinenumbers
	
	%%%%%%%%%% Discussion %%%%%%%%%%
\section{Discussion}
	\linenumbers
	% Overall summary of results and their meaning.
	
	
	
	% Issues with my model
		% the mass problem
			% Shrinking and how it addresses the growth speed issue in this model to a degree in my opinion.
		% gsi is likley underestimated in some cases at the very least 0.01 gsi wont be viable for all animals
	Some caveats with the results of the model.  First, growth is simulated as being extremely fast.  As previously stated, asymptotic size was reach by $\sim$10 days.  This is of course not representative of the real world, where individuals generally need several months to years to reach maturity.	
	The rapid growth may be due to several factors.  First is that metabolic cost may be underestimated.  Similar to \citeauthor{West2001} (\citeyear{West2001}), this study used resting metabolic rate to define metabolic costs.  However, this does not take other costs into account such as digestion and locomotion.
	%A rate which is typically used within energetic and growth studies, including early OGMs \parencite{Brett1971,Gillooly2001}  % ref for RMR in growth Tsoukali?  Burger? 
	This was addressed in traditional OGMs by \citeauthor{Hou2008} (\citeyear{Hou2008}).  However, due to the use of asymptotic mass in the parametrisation of this change, the same changes could not be used in this model.  Resting metabolic rate and active metabolic rates do not scale in the same way with mass \parencite{Gillooly2001, Weibel2004}.  The additional cost of active metabolic rate would cause a steeper scaling within the metabolic cost term, leading to more gradual growth (Fig. \ref{scaling_plot}). 
	As such, inclusion of active metabolic rates, while challenging to measure directly and implement, is needed.
	Additionally, there may be behavioural or physiological factors that would lead to an altered metabolic rate.  
	% TODO V presenting this as an issue but it is an area for improvement needs rephrasing and/or shifting
	In this regard, temperature plays a critical role.  It is well documented that a change in temperature will change many biological rates  (see \citeauthor{Peters1983} (\citeyear{Peters1983}), \citeauthor{Gillooly2001} (\citeyear{Gillooly2001}) etc.).  
	It has been shown that growth is dependent on temperature within fish. % sources
	For example, a 1.5\textdegree{}C increase in sea temperatures could result in a 15\% decrease in fish lengths \parencite{VanRijn2017}. % try to phrase better or bolster since it is lifted from diego
	The functional response data used in this study is standardised around 15\textdegree{}C \parencite{Pawar2012}.
	Meanwhile the metabolic cost is for an unspecified temperature \parencite{Peters1983}.  Using rates where the temperature effect is taken into account is crucial for model accuracy given the variation in rates that occurs over different temperatures.  Work such as \citeauthor{Barneche2014} (\citeyear{Barneche2014}) has investigated this effect, however the estimate for metabolic rate is several orders of magnitude lower than what was reported by \citeauthor{Peters1983} (\citeyear{Peters1983}), which is the rate used in this study.  Thus further investigation is required.  % should include a plot in the SI for this.
	
	In fish metabolic rate has also been shown to drop under starvation \parencite{Cook2000}.  In homeotherms feeding restriction has also been shown to also lower body temperature, since metabolic rate and core temperature are closely related in homeotherms \parencite{Ballor1991, Blanc2003,}.  
	
	Another possible point of error is the estimates for supply.  The parameters used from \citeauthor{Pawar2012} (\citeyear{Pawar2012}) are for a spectrum of animals from mammals to insects.  It is possible by reanalysing the data for only marine species, or more specifically only within taxon or species, predictions of supply could be improved \parencite{Marshall2019}.  % check dim paper for addressing rates in different environments
	
	A factor that is not taken into account in this model is that resources are not constant over time.  This can be implemented within the model by varying resource density over time.  The functional response will respond accordingly giving intake which varies through time.  One concern with implementing such a response is fluctuations are likely not experienced by all organisms in the same way.  For a fish with a small range a local fluctuation can be measured and described relatively simply.  However, for a fish with a very large range there is the possiblity leaving resource poor areas in search of richer waters.
	
	
	Despite these caveats, the patterns which arise from the model are promising.  The qualitative patterns seen in this study should not change even with slower growth.  This is because while the exact results are not representative of reality, the relationship between the values is.  For example, the manner in which supply and metabolic rate interact for different values of the metabolic scaling exponent.  The fact that a higher metabolic scaling exponent will cause intake and metabolic cost to intersect sooner does not change regardless of the absolute values.
	
 
	% However such field data is sparse. <- justify saying this some how. need a review or something.
	
	
	% What else can be done to address these issues
		% What is the minimum amount of GSI needed. reproductive structure will vary a bit between species and there is likely a minimum that a species must devote to these structures for them to be viable.
	
	
	% Emphasis that this is applicable to more than just fish - What needs to be done to apply the model to other animal groups.
	In conclusion, the model presented in this study is a promising base which can be expanded upon in a way that was not possible with previous OGMs allowing for much more controlled and detailed explanations of the factors controlling growth.  In contrast to all previous work, which assume optimal supply, the concept of varying supply is addressed using functional responses.  Additionally, qualitative evidence is provided supporting hyperallometric scaling in fish using energy budget as the basis.  The model can easily be applied to any animal taxon, not just fish, with some simple changes.  Additionally, there are clear directions to be explored to improve the model's accuracy.
	
	
	
	
	
	\begin{figure}
		\centering
		\includegraphics[width=\linewidth]{../../results/scalingplot.pdf}
		\caption{a) shows how maintenance cost outscales supply in a traditional OGM.  Growth only stops when maintenance reaches the supply line  b) shows scaling for supply and maintenance as equal. Since scaling is equal growth will never stop until the new cost of reproduction is introduced some time during development.  Blue line is supply, orange line is maintenance, green dotted line is reproduction}
		\label{scaling_plot}
	\end{figure}
%	\subsection{planned outline}
%	Overall summary of results and their meaning.
%	\\- I show that hyper allometry is possible within fish.  
%	\\- higher metabolic cost are leading to hyper allometry, even though I am seeing this for what the data considers too high (i.e. metabolic scaling expontent = 1) \parencite{Marshall2019} % see my google drive doc
%	\\- scaling in 3D will outscale metabolic cost in all cases Table \ref{parameters} because it is greater than 1, need reproductive cost to stop growth.
%	Shrinking and how it addresses the growth speed issue in this model to a degree in my opinion.
%	\\- Is like seeing the stage of growth we expect to see in a typical growth curve in reverse.  Where instead of growing into the asymptotic mass, it shrinks to it.  The energy balance remains the same.
%	\\
%	Issues with my model
%	- mass is too high, addressing a more detailed reproductive mass/metabolic rate may address this or reanalysing samraats data to be more taxon specific may do it.
%	\\
%	What else can be done.
%	\\- more fish specific intake using \cite{Pawar2012} methods as argued in \cite{Marshall2019}
%	\\- trying to model specific species with know growth curves to see how well our results map to them even if the curve isn't great it could be that reproductive output and final mass map well.
%	\\- more detailed metabolic cost since here is basically just RMR with higher exponents.  It may be that the constant shift too and this would lead to better predictions
%	\\- Incorporate oscillating resource density.  was attempted here but due to computational problems could not be implemented.  There is a lack of oscillation data this though and the goal would need to be far more specific since it will depend on scale and range of the organism.
%	\\- This same framework could be applied to other taxa with just a little adjustment (metabolic cost and reproductive equation?)
%	\\- There may be behavioural changes that occur at under feeding restriction that slows metabolism, similar to hibernation etc.
%	\\- temperature could be added here to adjust the metabolic and functional response rates %samraats data is standardised to 15C doesn't seem to be noted in peter's, hemingsen may show it
%	- develop the model to include storage like DEB or put this into a DEB framework.
%	\\
%	What is the minimum amount of GSI needed. reproductive structure will vary a bit between species and there is likely a minimum that a species must devote to these structures for them to be viable. %\cite{Gunderson1997} shows that mortality increases with higher GSI and larger inds have lower GSI does show low values for GSI (0.02) 
%	%GSI can reach 0.75 in invertebrates \cite{Parker2018}
%	\\
%	Emphasis that this is applicable to more than just fish - What needs to be done to apply the model to other animal groups. % need to open up gsi \cite{Parker2018} 0.75 inverts
%	\\
%	This is all energetically based, allowing us to ask questions of why or how there are exceptions to this in the case that organsims should be able to grow larger energetically speaking but may be limited genetically or via some other mechanism
	
	\nolinenumbers
%	\subsection{notes}
%	there are some integration issues that could allow for the appearance of more possible values
%	The shrinking helps to show what could be possible given there is a more gentle growth curve since the main difference is that the time when growth become more limited is just shown as shrinking vs slower growth.


	%%%%%%%%%% Conclusion %%%%%%%%%%
\section{Conclusion}
	\linenumbers
	
	
	
	\nolinenumbers
	
	%%%%%%%%%% Data Declaration %%%%%%%%%%
	\textbf{Code and Data Availability}

	Code is available at: \url{https://github.com/Don-Burns/Masters_Project}
	
	
	%%%%%%%%%% Bibliography %%%%%%%%%%
	\newpage
	\linenumbers
	
	%nicer indentation
	\addcontentsline{toc}{section}{\protect\numberline{}References}	
	
	\printbibliography
	
	% prob more official method of doing the above
%	\printbibliography[heading=bibintoc, title={References}]
	\nolinenumbers
	
	
	%%%%%%%%%% SI %%%%%%%%%%
	\newpage
	%to have table and figures numbered from 1 with prefix
\newcommand{\beginsupplement}{%
	\addcontentsline{toc}{section}{\protect\numberline{}Supplementary Information}
	\setcounter{table}{0}
	\renewcommand{\thetable}{S\arabic{table}}%
	\setcounter{figure}{0}
	\renewcommand{\thefigure}{S\arabic{figure}}%
}

\section*{Supplementary Information}
\beginsupplement

\subsection*{notes}
need section on value conversions and derivations

move any unreferenced sensitivity analyses here.


\begin{figure}
	\centering
	\includegraphics[width=\textwidth]{../../results/Sens_Maturation_Time_Long}
	\caption{Long time scale maturation time}
	\label{fig:sensmaturationtimelong}
\end{figure}

\begin{figure}
	\centering
	\includegraphics[width=\textwidth]{../../results/Sens_Maturation_Time_short}
	\caption{Short time scale maturation time}
	\label{fig:sensmaturationtimeshort}
\end{figure}

\begin{figure}
	\includegraphics[width=\textwidth]{../../results/Sens_Low_Resources_Maturation_Time_short}
	\caption{Low Resources Maturation}
	\label{fig:senslowresourcesmaturationtimeshort}
\end{figure}

\begin{figure}
	\centering
	\includegraphics[width=\textwidth]{../../results/Sens_Metabolic_Exponent}
	\caption{High Resources Metabolic Exponent}
	\label{fig:sensmetabolicexponent}
\end{figure}

\begin{figure}
	\centering
	\includegraphics[width=\textwidth]{../../results/Sens_Low_Resources_Metabolic_Exponent}
	\caption{Low Resource Metabolic Exponent}
	\label{fig:senslowresourcesmetabolicexponent}
\end{figure}

\begin{figure}
	\centering
	\includegraphics[width=\textwidth]{../../results/Sens_Resource_Density}
	\caption{Large scale Resource Resource Density}
	\label{fig:sensresourcedensity}
\end{figure}

\begin{figure}
	\centering
	\includegraphics[width=\textwidth]{../../results/Sens_Low_Resources_Resource_Density}
	\caption{Fine scale resource Density}
	\label{fig:senslowresourcesresourcedensity}
\end{figure}

\begin{figure}
	\centering
	\includegraphics[width=\textwidth]{../../results/Sens_Shrink}
	\caption{High Resources Shrinking}
	\label{fig:sensshrink}
\end{figure}

\begin{figure}
	\centering
	\includegraphics[width=\textwidth]{../../results/Sens_Low_Resources_Shrink}
	\caption{Low resources shrinking, missing values around 0 are because no viable solutions could be found.}
	\label{fig:senslowresourcesshrink}
\end{figure}




	


	
	
\end{document}} % to include word count 

%%%%%%%%%% Formatting %%%%%%%%%%
\usepackage[margin=2cm]{geometry} % margins of 2cm
\linespread{1.5} %1.5 spacing
%\renewcommand{\familydefault}{\sfdefault} % set font to arial clone (helvet)

\usepackage[compact]{titlesec} % reduce spacing bewteen section titles

%%%%%%%%%% Bibliography %%%%%%%%%%
\addbibresource{../Masters_Thesis.bib}

%commands 

%%%%%%%%%% Document %%%%%%%%%%
\begin{document}
	%%%%%%%%%% Title Page %%%%%%%%%%
	%\newcommand{\crest}{\includegraphics[width = 4cm, keepaspectratio]{../images/IC_Crest.eps}} % Imperial crest
%%formating
\begin{titlepage} % Suppresses headers and footers on the title page
	
	\centering % Centre everything on the title page
	
	\scshape % Use small caps for all text on the title page
	
%	\vspace*{\baselineskip} % White space at the top of the page
	
	%------------------------------------------------
	%	Title
	%------------------------------------------------
	
	\rule{\textwidth}{1.6pt}\vspace*{-\baselineskip}\vspace*{2pt} % Thick horizontal rule
	\rule{\textwidth}{0.4pt} % Thin horizontal rule
	
	\vspace{0.75\baselineskip} % Whitespace above the title
	
	{\LARGE Using Realistic Resource Supply to Show Hyperallometric Reproductive Output in Fish\\} % Title
	
	\vspace{0.75\baselineskip} % Whitespace below the title
	
	\rule{\textwidth}{0.4pt}\vspace*{-\baselineskip}\vspace{3.2pt} % Thin horizontal rule
	\rule{\textwidth}{1.6pt} % Thick horizontal rule
	
	\vspace{1\baselineskip} % Whitespace after the title block
	
	%------------------------------------------------
	%	Subtitle
	%------------------------------------------------
	
	%SUBTITLE? % Subtitle or further description
%	Student:
	
	
	\vspace{0.5\baselineskip} % Whitespace before 
	
	{\scshape\Large D\'onal Burns  \\} % my name
	
	\vspace{0.5\baselineskip} % Whitespace below 
	
	\textit{CID: 01749638 \\ Imperial College London \\ Email: donal.burns@imperial.ac.uk} % affiliation and email
	
	\vspace*{2\baselineskip} % Whitespace under the subtitle
	
	
	
%	Supervisor:
%	
%	
%	\vspace{0.5\baselineskip} % Whitespace before 
%	
%	{\scshape\Large Samraat Pawar \\} % supervisor name
%	
%	\vspace{0.5\baselineskip} % Whitespace below 
%	
%	\textit{Imperial College London \\ Email: s.pawar@imperial.ac.uk} % affiliation and email
%	
%	\vspace{3cm} % Whitespace between 
	

	
	%% crest
	
	\includegraphics[width = 4cm, keepaspectratio]{../images/IC_crest.pdf}
	
	\vspace{0.3\baselineskip} % Whitespace under the Uni logo
	
	Submitted: August 27$^{th}$ 2020 % Publication Date
	
	

		%%Submission clause
	\vspace{0.5\baselineskip}
	
	A thesis submitted in partial fulfilment of the requirements for the degree of
	%Computational Methods in Ecology and Evolution 
	Master of Science at Imperial College London
	\vspace{0.5\baselineskip}
	
	Formatted in the journal style of Functional Ecology	
	\vspace{0.5\baselineskip}
	
	Submitted for the MSc in Computational Methods in Ecology and Evolution
	
\end{titlepage}

	

	%%%%%%%%%% Declaration %%%%%%%%%%
	\section*{Declaration}
	I declare this project as my own work.  The model presented here was developed in conjunction with my supervisor, Dr. Samraat Pawar, and Ph.D. students Tom Clegg and Olivia Morris.  I was responsible for any simulations and data presentation.\newline
	%% word count
	\textbf{Word Count: \wordcount}

	\newpage
	
	%%%%%%%%%% Abstract %%%%%%%%%%
	\section*{Abstract}
	\linenumbers
%	Size is essential to reproductive output. By extension understanding growth, determines size allows understanding of reproductive output. 
	 
	With recent results showing that reproduction in fish scales hyperallometrically there is a need to update growth OGMs to reflect this fact.  Current OGMs assume optimal intake, an assumption which is not always reflected in the field.  In this study I develop an energy intake focused approach to explaining growth, an area which has not been covered within current literature, and shows that hyperallometric scaling of reproductive output arises when allowing for variable reproductive scaling and maximising for fitness.  The model is applicable to not only fish, but any animals taxon with some simple parameter adjustments.  I offer direction for improvements and areas to be developed in order to allow the model to be applicable to any temperature range.
	\vspace*{0.5 cm}
	\newline
%	TODO pieces to include:
%	
%	show the effects of the exponent and why it is important that we use the correct one moving forward.
%	give direction as to what needs to be done going forward.  what areas need work
	%TODO KEywords
	\textbf{Keywords:}\\
	allometry; functional response; metabolic theory; growth; intake; life history; metabolism; reproduction; reproductive output; supply

	
	
	\nolinenumbers
	%%%%%%%%%% Acknowledgements %%%%%%%%%%
	%\thispagestyle{empty}

\mbox{}\newline\vspace{10mm} \mbox{}\LARGE
%
{\bf Acknowledgements} \normalsize \vspace{5mm}

%I would like to thank my supervisor.....
%thanks to luke who gave a great base to start from on the project




 %TODO order with abstract since not line numbered?
	
	%%%%%%%%%% Table of Contents %%%%%%%%%%
	\tableofcontents
	\newpage

	%%%%%%%%%% Introduction %%%%%%%%%%
\section{Introduction}
	\linenumbers
%	TODO Include Diego's references

	%Ease into it a bit first
	Body mass plays a major role in determining many biological factors.  For example, larger individuals are less vulnerable to predation, have lower mass specific metabolic rates and produce more offspring in their lifetime \parencite{Peters1983, Barneche2018, Craig2006, Magnhagen2001, Hixon2014, Marshall2006}.
	% growth, what we do and don't know.
	By extension, knowing how body mass changes over an organisms lifetime is the gateway to understanding how many biological rates change throughout ontogeny.  This is because so many biological rates scale with mass \parencite{Kleiber1932}.  Despite its importance, relatively little is know about the factors determining growth trajectories \parencite{Arendt2011, Marshall2019}.
	\\
	% TODO bring in OGM growth curve and use to illustate the basic OGM as per samraats comments
	% Why is growth important
	In the case of fish, understanding growth and what factors play a role in determining it is not only insightful from the perspective of understanding the world around us, but can also be used to better manage the many fisheries and marine protected areas around the world \parencite{Heino2013, Lester2009}.  An objective which is becoming increasingly important as the oceans' fish stocks continue to be depleted by overfishing. 
	%VV rephrase this part to hint more towards size than it currently does VV
	This is compounded by global warming which threatens to alter the structure of the ecosystems even if they are not fished and left in their ``natural" state \parencite{Bruno2018}.
	It is already known that metabolic rate is dependant on temperature which in turn will lower fish sizes \parencite{Gillooly2001, Brown2004} and with increasing global temperatures, understanding in greater detail how increased metabolic rates %TODO mentioning metabolism here is now a bit random that it has been moved
	may affect growth is useful in population management.
	\\
	% introduction of models
	To date many models have been developed to predict and describe the growth of an organism throughout its lifetime.  The three main approaches used are the von Bertalanffy model, the dynamic energy budget (DEB) model and the ontogenetic growth model (OGM), which will be the focus of this study \parencite{Putter1918, vonBertalanffy1938, Kooijman1986, West2001}.  All of these are energetic based models with varying assumptions, key among which is the scaling of resource supply and metabolic rate with mass. %TODO move this final sentence to before the previous so the para end with this study focusing on ogm methodology  
	%In OGMs supply is thought of as being optimal at all times which leads to the assumption that intake scales with mass to the power of 0.75.  Indeed while under optimal conditions this may be true, it neglects that this situation is thought to rarely occur in the field \parencite{Pawar2012}.
	%	introduce \cite{West2001}	
	\\
	One of the best known examples of an OGM is the model developed by \cite{West2001}.  This model is parameterised around the average energy content of animal tissue and asymptotic mass.  Asymptotic mass being the mass at which growth has essentially stopped due to metabolic cost and energy intake equalling each other (Fig. \ref{scaling_plot}a). The model hinges on the scaling with mass between energy intake (m$^{0.75}$, allometric sub-linear scaling) and maintenance cost (m$^1$, isometric linear scaling).  In other words, as mass increases maintenance costs will slowly overtake the intake rate and halt growth (Fig. \ref{scaling_plot} a).  	
	% talk about determinant and indeterminant growth/
	% move to including reproduction with \cite{Charnov2001} mashed with \cite{Hou2008} imporvements briefly
	The framework used by \citeauthor{West2001} (\citeyear{West2001}) was latter developed by \citeauthor{Charnov2001} (\citeyear{Charnov2001}) to take the cost of reproduction into account and allowing the estimation of lifetime production of offspring.  \citeauthor{Hou2008} (\citeyear{Hou2008})  developed \citeauthor{West2001}'s model further by expanding maintenance cost to include the cost of feeding and digestion (specific dynamic action), synthesis of new tissue and activity.
	% begin caveats
	% tautology still present 
	% Discuss allometry and isometry here to highlight what scaling super or sub linearly means
	In the above OGMs intake is assumed to scale sub-linearly to the power of 0.75.  This is due to the assumption that the individuals are consuming at an optimal rate at all times and therefore the only limitation is their ability to make use of that energy.  In this case, intake should theoretically scale to the power of 0.75 (see \cite{West1997}).  However, this is not always the case in the field.  It has been shown that, for non-optimal consumption, steeper scaling can occur \parencite{Pawar2012, Peters1983}. Additionally, OGMs, like many growth and metabolic models, typically use basal or resting metabolic rate to calculate metabolic cost.  This is the minimal metabolic rate of an organism and is typically thought of as the rate of the organism when relaxed and at rest.  However, it has been shown, once factors such as movement are taken into account, that the scaling becomes steeper \parencite{Weibel2004}.
	\\
	% tautology of OGMs and \cite{Hou2011} close but still issues
	The issue of non-optimal feeding is addressed somewhat by \citeauthor{Hou2011} (\citeyear{Hou2011}).  However, this growth was only investigated as, essentially, a proportion of optimal consumption and does not address a potential change in scaling of intake rate.
	Another limitation of the models used in previous OGMs is the dependence on asymptotic mass.  The models are entirely dependent on the value of optimal intake and asymptotic mass.  All other values, such as metabolic cost, are then derived in relation to these.  However, organisms are not born with some inherent restriction to the size they can attain, at least not energetically.  If there is surplus energy for a given mass the organism should be able to grow.  Relying on asymptotic mass to define the upper bound of attainable mass does not allow for investigation of the mechanisms that underpin asymptotic size in reality. 
	\begin{figure}[h]
		\centering
		\includegraphics[width=\linewidth]{../../results/scalingplot.pdf}
		\caption{a) shows how maintenance cost out-scales supply in a traditional OGM.  Growth only stops when maintenance (scaling exponent = 1) reaches the supply line (scaling exponent = 0.75)  b) shows scaling for supply and maintenance as equal. Since scaling is equal growth will never stop until the new cost of reproduction is introduced some time during development.  Blue line is supply, orange line is maintenance, green dotted line is reproduction}
		\label{scaling_plot}
	\end{figure}
	With two of the key assumptions of current OGMs, that reproduction and metabolism scale isometrically, not holding in the field \parencite{Barneche2018, Pawar2012, Peters1983} there is a need to take an unexplored approach to modelling fish growth, in particular choosing to focus on developing how intake is described so as to better reflect the real world.  To do this an obvious starting point is to model intake as a functional response \parencite{Holling1959} so as to better reflect real world intake rates in terms of consumed biomass over time.  Non-optimal supply is a currently unexplored area within growth modelling.  This is likely due to the difficulty of directly measuring intake, especially in the field, Perhaps as a result comparatively less is known about consumption.  This leads to the need for proxy values to estimate it, for example nutrient flux \parencite{Schiettekatte2020}, or drawing broad relationships to approximate consumption, as this study will do.
	Changing the manner in which intake is defined also requires changing metabolic cost, because the two are dependent upon each other in current OGMs.  This can be done by defining metabolic rate as a value dependent on current mass rather than asymptotic mass as has been done in OGMs up until this point.  This thought process is more mechanistic since an organism has no concept of ``How large should I grow?", but rather will acquire as much resources as it capable of at its current life stage and size.  Taking this more bottom-up mechanistic approach also allows exploration of what controls how large fish grow to, since as previously mentioned, from an energetic standpoint, an organism can grow indefinitely as long as there is surplus energy available after costs have been paid.  Of course, there are also mechanical and genetic limitations upon organism size, but once size is constrained to what is known to exist, this should not be an issue.  %TODO fix "this should not be an issue"
	\\
	This study takes the novel approach of using a mass-specific functional response and assimilation efficiency to describe how intake changes both throughout ontogeny and varying levels of resource availability. This study focuses on supply and growth within fish, however the same principles can be applied to other taxa.
	\\
	Previous OGMs have assumed that reproduction scales isometrically with mass.  This is indeed the case, within fish larger individuals produce more offspring than smaller ones.  % this is known in general isometrically
	However, it has been shown that larger fish produce far more offspring than the equivalent mass composed of smaller fish,  that is a 2kg fish will produce more offspring than two 1kg fish, i.e. reproduction scales hyperallometrically \parencite{Barneche2018}.
	Larger fish also use energy more efficiently than multiple smaller ones per unit mass.  This is due to them having a lower mass specific metabolic rate \parencite{Peters1983, Kleiber1932, Brown2004}.  
	% larger mothers produce larger offspring which may better survive \cite{Barneche2018}
	Additionally, larger mothers produce larger offspring, which are then more likely to survive to adulthood and reproduce \parencite{Hixon2014, Marshall2006}. 
	This has led to thinking rather than metabolism having steeper scaling than supply being the reason that growth stops (Fig. \ref{scaling_plot}a).  Instead, metabolism and supply at saturated resources should scale in a similar manner, with the trigger for growth slowing and stopping being the scaling of reproduction (Fig. \ref{scaling_plot}b) \parencite{Marshall2019, Sibly2020}.
	\\
	% Justify and appeal my methods	
	Assuming that fish have evolved to maximise reproductive output and can adapt to find an optimal strategy within the constraints of resource density, simulations can be carried out to demonstrate what conditions need to be met in order to achieve hyperallometric scaling of reproduction from an energetic perspective.  This study will answer 1) under what metabolic scaling exponent can hyperallometric reproduction arise, %TODO this may be a bit of an aggressive statement, tone down? Show that metabolic exp plays a large role in rho or something like that.
	2) possible scaling of metabolism and reproduction is dependent upon supply and by extension dimensionality.

	\nolinenumbers
	
	%%%%%%%%%% Methods %%%%%%%%%%
\section{Methods}
	\linenumbers
	
	\subsection{Altering OGMs to Account for Resource Supply}
	In order to address the issue of supply in the context of an OGM, which can be generically described as $dm/dt = gain - loss$, some changes need to be made to the model's terms.  The first is to remove the assumption of asymptotic mass and the reliance of cost upon it.  Within a tradional OGM the gain term ($ a $) and asymptotic mass are used to define the metabolic cost ($ b $).  However, since the assumption of perfect intake is going to be broken, because of the variable supply, this relationship no longer holds.  As such, both intake and metabolic cost need to be redefined.  Additionally, in light of recent work showing that reproduction scales allometrically and not isometrically, the reproductive cost must also be modified from the form used by \citeauthor{Charnov2001} (\citeyear{Charnov2001}) \parencite{Marshall2019, Barneche2018} %TODO check marshal ref data origin
	
	
	\subsubsection{Gain}
	To define intake a natural starting place is the functional response \parencite{Holling1959}.  Functional responses  are used to define how much an organism consumes for a given resource density and is described by the following equation:	
	\begin{equation}
		\label{functional_repsonse}
		f(\cdot) = \frac{a X_r}{1 + a h X_r}
	\end{equation}
	where, $ f(\cdot) $ is the functional response, $ a $ is the search rate, $ h $ is handling time and $ X_r $ is resource density.  
	For a fixed mass and increasing resource density Eq \ref{functional_repsonse} produces a sigmoidal shape with intake eventually reaching an asymptote after some saturating amount of resources is reached.  At lower resource densities, the intake rate is primarily defined by the search rate with higher search rates yielding higher intake rates.  Conversely, at high resource densities, intake rate is approximately equal to the inverse of the handling time ($ h^{-1} $), where lower handling times yield higher intake rates.  
	
	An organism's functional response will not remain constant throughout its life history.  Search rate and handling time are affected by both the organism's mass and how it interacts with its environment \parencite{Pawar2012}.  
%	Within this model mass will be known for all time points since that is one of the quantities being predicted.  
	Interactions can be broken into 3D and 2D, that is whether the organism consumes from a 2D ``surface" e.g. a cow grazing or a 3D ``volume" e.g. a pelagic consumer which consumes prey from within the water column.  As such, both handling time and search rate can be defined as Eq. \ref{handling_time} and Eq. \ref{search_rate} respectively.
	\begin{equation}
		\label{search_rate}
		a(m) = a_0 m_t^\gamma
	\end{equation}
	
	\begin{equation}
		\label{handling_time}
		h(m) = t_{h,0} m_t^\beta
	\end{equation}
	A functional response alone is not enough to fully define intake.  This is because processing of consumed resources is not one hundred percent efficient which leads to inevitable loss consumed energy.  As a result, to achieve the final gain term, a dimensionless efficiency term $\epsilon$ is applied.  In poikilotherms assimilation efficiency is roughly 70\% \parencite{Peters1983}
	
	\subsubsection{Loss}
	Metabolic cost  has previously been dependant upon the gain term within traditional OGMs (see \cite{West2001, Hou2008}).  However, for non-maximal intake the relationship will no longer hold true.  As a result, this model has taken previously measured values to be used as metabolic cost (see Eq. \ref{metabolic_cost} taken from \cite{Peters1983} % add ref to hemmingsen here since that is the peters source
	and Table \ref{parameters} for further details).
	\begin{equation}
		\label{metabolic_cost}
		B_m = 0.14 m_t^\mu
	\end{equation}
	Next to take allometric scaling of reproduction into account, the reproductive cost term from \citeauthor{Charnov2001} (\citeyear{Charnov2001}) is changed from $cm^1$ which assumes isometric scaling to $cm^\rho$.  $c$ can be interpreted as the proportion of mass dedicated to reproduction, i.e. the gonadosomatic index (GSI) of the fish \parencite{Charnov2001}.  Just as in \citeauthor{Charnov2001} (\citeyear{Charnov2001}) reproductive cost is only taken into account once maturity is reached.  Meaning that until a length of time ($\alpha$) has passed, reproductive cost is considered to be zero.
	
	\subsubsection{Full Growth Equation}
	When all newly defined terms are taken together they produce Eq. \ref{dmdt_juvenile}  and \ref{dmdt_mature} which describe the change in mass for a fish before (Eq. \ref{dmdt_juvenile}) and after maturation (Eq. \ref{dmdt_mature}).
	\begin{align}
		\label{dmdt_juvenile}
		\frac{dm}{dt} &= \epsilon f(\cdot) - B_m & m_t < m_\alpha \\
		\label{dmdt_mature}
		\frac{dm}{dt} &= \epsilon f(\cdot) - B_m - cm_t^\rho & m_t \geq m_\alpha
	\end{align}

	
	
	\subsection{Calculating Fitness}
	At any time ($ t $) a reproducing organism devotes some amount of energy to reproduction.  This is the product between the amount of mass dedicated to reproduction ($ cm^\rho $) and a declining efficiency term ($ h_t $) which begins at maturity ($ \alpha $) and represents reproductive senescence \parencite{Benoit2018, Vrtilek2018, Stearns2000}.  In addition to amount of reproduction, the offspring are also subject to mortality ($ l_t $).  By combining the two, lifetime reproductive output can be estimated and is described by the ``characteristic equation" (Eq. \ref{characteristic_equation}) which represents reproductive output in a non-growing population \parencite{Tsoukali2016, roff1993, Roff2001, stearns1992evolution, Arendt2011, Roff1986, Roff1984}
	\begin{equation}
		\label{characteristic_equation}
		R_0 = \int c m_t^\rho h_t l_t 
	\end{equation}
	Mortality is experienced differently by juvenile ($ t < \alpha $) and reproducing individuals ($ t \leq \alpha $). % TODO <-dont like this sentence
	Mortality of offspring prior to maturity is described as a survival rate $ l_t = e^{-Z(t)} $ which is an exponentially decreasing function bounded between zero and one.  It controls how many offspring make it to maturity.  After maturity, survival is again described as an exponential function which takes time to maturity into account, $ l_t = e^{-Z(t-\alpha)} $.  
	Reproductive senescence can be also be estimated as an exponential function which begins after maturity and declines over time  ($ e^{-k(t-\alpha)} $), where $ k $ is the senescence term.  When all values are inserted into the characteristic equation, it results in the equation used by \citeauthor{Charnov2001} (\citeyear{Charnov2001}) with the inclusion of reproductive senescence (Eq. \ref{reproductive_output}).
	\begin{equation}
		\label{reproductive_output}
		R_0 = c\int_0^\alpha e^{-Z_t} dt  \int_\alpha^\infty m_t^\rho e^{-(\kappa + Z)(t - \alpha)}dt\\
	\end{equation} 
	In Eq. \ref{reproductive_output}, $ Z $ represents instantaneous mortality.  This rate has been shown to be related to time of maturation in many taxon groups, and within it follows the relationship $ \alpha \cdot Z \approx  2$.  This can then be rearranged to estimate instantaneous mortality, $ Z \approx 2/\alpha  $
	
	\subsubsection{Maximising Reproduction}
	It is assumed that evolution will converge on metabolic values which maximise fitness. 
	Fitness being defined as how much an individual is able to contribute to the gene pool \parencite{Speakman2008, Stearns2000}.  % Would like to remove this and link more smoothly
	To this end, lifetime reproductive output is often used as a measure of fitness \parencite{Charnov1991, Audzijonyte2018, Speakman2008, Stearns2000, Charnov2001, Tsoukali2016, Brown1993, Charnov2007}.  Therefore, by maximising for reproductive output, it should become clear what parameters will yield the highest fitness.  These parameters will then show whether, within a theoretical framework, hyperallometric scaling arises.
	
	To find all optimal values for reproduction would require Eq. \ref{reproductive_output} to be solved analytically.  However, since no such solution is possible, I simulated the problem numerically to obtain a result.  This was done by simulating across values of $ c $ and $ \rho $, the parameters of interest between growth (Eq. \ref{dmdt_juvenile} and \ref{dmdt_mature}) and reproductive output (Eq. \ref{reproductive_output}).  $ c $ was bound between 0 and 0.4, which encapsulates the values measured within fish \parencite{Benoit2018, Roff1983, Fontoura2009}.  Though it has been shown to reach as much as 0.7 in some species \parencite{Parker2018}.  To search for any hyperallometry within reproduction, $ \rho $ was bound between 0 and 2.  
	The simulation was then run at 0.01 value intervals in both $c$ and $\rho$ over a lifespan of one million days.  The results of each simulation were recorded and any non-viable results were discarded.  A result was considered non-viable if fish had ``shrunk" more than 5\% in order to accommodate reproductive costs.  Shrinking occurs in the model because  the combined loss of energy to metabolism and reproduction is too much for the simulated values at the mass achieved by maturation so the individual experiences a deficit of energy, which is paid by loss in mass until equilibrium is achieved. % \cite{VandenBerghe1992} for reproductive mass loss (though it is due to behaviour changes)
	Shrinking is not expected at maturity in reality.  Typically, maturity will occur while the organism still has room for growth.  It is the onset of reproduction which is considered to slow or stop growth % this is the case where the metabolic exponent is the same or less than the intake one.
	 (see Fig. \ref{OGM_Curve}).  Shrinking can be thought of as starvation in a real organism.  If energetic cost are not met then energy reserves in the body, such as fat and muscle, are broken down for energy.  It has been shown that some fish are capable of losing up to 10\% of their body mass \parencite{VandenBerghe1992}.  However, this was during the breeding season and caused by behavioural changes due to parenting.  Additionally, individuals were shown to rebound back to their``normal" body mass once the breeding season had ended.
	% survival impacts of shrinking? are there sources?
	% \cite{Lambert2000}
	
	\subsection{Sensitivity Analysis}

	In order to determine the roles of metabolic exponent, maturation time and resource density within the model, sensitivity analyses were performed on each parameter with regard to $c$ and $\rho$.  This was done by simulating the parameters across multiple values and obtaining the optimal value for $c$ and $\rho$ as described above.
	
	
	\begin{centering}
		
		%		TODO check description capitalisation
		
		\begin{table}[h!]
			
			\caption{Table describing parameters used in the model, along with values, units and sources where applicable.  The units of resource density change depending on the dimension of intake.  $m^D$ represent either $m^2$ in 2D or $m^3$ in 3D} 
			\label{parameters}
			\vspace{2mm}
			\begin{tabular}{c p{3cm} l l l p{3cm}}
				\hline
				Parameter 	& Description 			& Value 	& Units 	& Range 		& Source \\
				\hline
				$m$			& Mass					& -			& kg day$^{-1}$& -			&		\\
				
				$B_m$		& Metabolic Cost		& $0.14 m^{\mu}$ & kg day$^{-1}$& - 	& \cite{Peters1983}\\
				$\mu$		& Metabolic Exponent	& -			&	-		& 0.75 - 1.0	& - \\
				$\alpha$	& Age of maturity		& -     	& day		& -				& -\\
				$c$			& Reproduction scaling constant & - & kg day$^{-1}$& 0 - 0.5 		& -\\
				$\rho$		& Reproduction scaling exponent	& -	&	-		& 0 - 1.5			& -\\
				$Z$			& Rate of instantaneous mortality& $2/\alpha$	& & & \cite{Charnov2001}\\%double check ref
				$k$			& Reproductive senescence & 0.01\\
				
				$\epsilon$	& Assimilation Efficiency & 0.70 & - & - 		& \cite{Peters1983} \\
				$X_r$ 		& Resource Density		& -		& kg/m$^D$		& 0.11 - 30				& -\\
				$\gamma$	& Search rate scaling exponent & 0.68 (2D)	& - & - & \cite{Pawar2012} \\
				&						& 1.05 (3D)\\
				$a_0$		& Search rate scaling constant & $10^{-3.08}$ (2D) & m$^2$ s$^{-1}$ kg$^{-0.68}$   & - &\cite{Pawar2012}	\\
				&						& $10^{-1.77}$ (3D)& m$^2$ s$^{-1}$ kg$^{-1.05} $\\
				$\beta$		& Handling time scaling exponent& 0.75 & - & - & \cite{Pawar2012}\\
				$t_{h, 0}$	& Handling time scaling constant& $10^{3.95}$ (2D) &kg$^{1-\beta}$ s& -& \cite{Pawar2012}	\\
				&						& $10^{3.04}$ (3D)			&kg$^{1-\beta}$ s\\
				\hline
			\end{tabular}
		\end{table}
	\end{centering}

	\newpage

	\nolinenumbers
	%%%%%%%%%% Results %%%%%%%%%%
\section{Results}
	\linenumbers
	
%	\subsection{notes}
%		scenesence not much of a factor
% 10 \% predicted under Roff 1983

	\subsection{Growth Curve and Maturation Time}
	% talking about at meta_exp = 1
	As can be seen in Figure \ref{growth_curve}, growth is very fast within the model.  Asymptotic mass is reached by $\sim$15 days.  This makes interpreting any results regarding maturation time difficult % TODO < rephrase
	because any time after $\sim$15 days produces the same result (see Fig. \ref{fig:senslowresourcesmaturationtimeshortexp1} - \ref{fig:senshighresourcesmaturationtimeshortexp1}).
	However, as can be seen in Fig. \ref{fig:sensmaturationtimeshortexp075}, where maturation is occurring early during the growth phase, hyperallotric scaling emerges.  This pattern shows in particular at lower resource densities in 3D (Fig. \ref{fig:sensverylowresourcesmaturationtimeshortexp075} and \ref{fig:sensverylowresourcesmaturationtimeshortexp1})
	
	
	\subsection{Sensitivity Analysis}
	\subsubsection{Resource Density}
	The scaling relationship emerges as would be expected from the scaling of the functional response.  At low resource densities the output of the functional response will scale similarly to search rate.  The scaling of which is higher in 3D as resources increase and the response shifts to scaling similarly to the inverse of handling time.  At this point $\rho$ starts to take values which are higher in 2D than 3D, because of the higher normalisation constant in 2D.
	
	\subsubsection{Metabolic Exponent}
	% talking about at meta_exp = 1
	The expected result is for increasing values of metabolic scaling exponent ($\mu$) for $\rho$ to also increase.  This is because the lower values of $\mu$ will result in there being a larger gap between the scaling of intake and maintenance, which allows for steeper scaling in reproduction (see Fig. \ref{scaling_plot} b) \parencite{Marshall2019}.  This appears to be the case when analysing $\rho$ with respect to the other parameters where $\mu = 1$ or 0.75 (e.g. Fig. \ref{fig:senshighresourcesmaturationtimeshortexp1} and \ref{fig:sensmaturationtimeshortexp075} or Fig. \ref{fig:sensresourcedensityexp1broad} and \ref{fig:sensresourcedensityexp075fine}).  However, when $\mu$ is explicitly simulated over a variety of values, the trend suggests that increasing $\mu$ allows for higher values of $\rho$.  
	
	\subsubsection{Shrinking}
	% talking about at meta_exp = 1
	At saturated resource densities, allowing for greater proportions of shrinking enables $\rho$ to be set to larger values in both dimensions.  This is perhaps not surprising as increasing the shrinking proportion allows for a larger reproductive cost since all individuals will already be at asymptotic mass at time of maturation.
	
	\subsubsection{$c$ Values}
	Estimations of $c$ are low in many cases, especially in 2D.  While this may be low compared to the $\sim$10\% - 35\% expected \parencite{Benoit2018, Fontoura2009, Roff1983} it is not unprecedented for values of 2\% to be observed in some fish \parencite{Gunderson1997}.  It may be necessary for the lower bounds of $c$ to be adjusted based on what is expected or even viable in the organisms being simulated.
	
%	minimum viable resource is dependent upon meta.exp, with 0.75 allowing for less than that of 1 and scaling in between 0.2 for 1 and 0.15 for 0.75. likewise 3D is capable of going lower than 2D due to the higher intake rate.  shrinkage plays no role here.

	\begin{figure}[h!]
				
%		TODO 1.The title for each plot needs to be cleaned up.  2.The units should be m^2 and 3.need letter labels for each.
		\includegraphics[width=\textwidth]{../../results/report_3D2D_HighLowResCLEAN_meta_exp_1.pdf}
		
		\caption{Multiplot with 2D vs 3D and varying resource densities. The optimal value is denoted by the blue circle.
		The resources used for the low resource scenario (top row) is the minimum amount of resources that allows growth with a $c$ and $\rho$ of 0.01.  As would be expected, since has steeper scaling, 3D allows for growth at smaller resource densities than 2D.
		Low reources in 2D were $ \approx 0.1kg/m^2 $ and $ 0.00035kg/m^2 $ in 3D.
		$ 100kg/m^2 $ was used for high resources because, which although unrealistic, ensures that resources are not a limiting factor in the simulations.
		Hyperallometric scaling is observed in 2D ($\rho = 1.64$ at high resources and 1.99 at low resources (NOTE: This is because 2 is the upper limit of $\rho$ I have simulated here). 
		Scaling in 3D is slightly hypoallometric $rho = 0.94$ and c = 0.01 at high resources. $\rho = 0.8$ and $c = 0.01$ at low resources
		The metabolic scaling exponent = 1 in all cases}
		\label{resources2D3D_meta_exp1}
	\end{figure}

	\begin{figure}[h!]
		
		%		TODO 1.The title for each plot needs to be cleaned up.  2.The units should be m^2 and 3.need letter labels for each. CHECK SHRINKING
		\includegraphics[width=\textwidth]{../../results/report_3D2D_HighLowResCLEAN_meta_exp_075.pdf}
		
		\caption{Multiplot with 2D vs 3D and varying resource densities. The optimal value is denoted by the blue circle.
			The resources used for the low resource scenario (top row) is the minimum amount of resources that allows growth with a $c$ and $\rho$ of 0.01.  As would be expected, since has steeper scaling, 3D allows for growth at smaller resource densities than 2D.
			Low reources in 2D were $ \approx 0.1kg/m^2 $ and $ 0.00035kg/m^2 $ in 3D.
			$ 100kg/m^2 $ was used for high resources because, which although unrealistic, ensures that resources are not a limiting factor in the simulations.
			Hyperallometric scaling is observed in 2D ($\rho = 1.64$ at high resources and 1.99 at low resources (NOTE: This is because 2 is the upper limit of $\rho$ I have simulated here). 
			Scaling in 3D is slightly hypoallometric $rho = 0.94$ and c = 0.01 at high resources. $\rho = 0.8$ and $c = 0.01$ at low resources
			The metabolic scaling exponent = 0.75 in all cases}
		\label{resources2D3D_meta_exp0.75}
	\end{figure}
	
%	\begin{figure}[h!]
%		
%		\includegraphics[width=\textwidth]{../../results/report_2D_exponents.pdf}
%		
%		\caption{Multiplot with changing metabolic exponents in 2D, ranging from 0.75 to 1, resource density = 0.01, time to maturation is 5 years.  
%		It can be seen that higher metabolic responses limit the proportion of mass that can be dedicated to reproduction.  However, in exchange, they allow for much steeper scaling.  A patter which is mirrored in 2D and 3D (fig. \ref{exponents3D})}
%		\label{exponents2D}
%	\end{figure}
%		
%	\begin{figure}[h!]
%		
%		\includegraphics[width=\textwidth]{../../results/report_3D_exponents.pdf}
%		
%		\caption{Multiplot with changing exponents in 3D, ranging from 0.75 to 1, resource density = 0.01, alpha is 5 years. See fig \ref{exponents2D} for explanation. (NOTE: either 2D or 3D likely belongs in SI) } %either of these likely belong in SI.
%		\label{exponents3D}
%	\end{figure}
		
	\begin{figure}
		\centering %TODO Units for axis?
		\includegraphics[width=\textwidth]{../../results/pretty_curve}
		\caption{Example of the growth curve and cumulative reproduction expected from a traditional OGM model. Maturation occurs at 1000 days, after which growth is less steep until reaching asymptotic mass.  }
		\label{OGM_Curve}
	\end{figure}
	
	\begin{figure}
		
		\includegraphics[width=\textwidth]{../../results/report_growth_curve.pdf}
		\caption{The growth over a fish which consumes in 2D.  Maturation occurs at 5 years (1825 days).  The fish was allowed to shrink by 5\% at the onset of reproduction.}
		\label{growth_curve}
	\end{figure}



	\nolinenumbers
	
	%%%%%%%%%% Discussion %%%%%%%%%%
\section{Discussion}
	\linenumbers
	% Overall summary of results and their meaning.
	
	
	
	% Issues with my model
		% the mass problem
			% Shrinking and how it addresses the growth speed issue in this model to a degree in my opinion.
		% gsi is likley underestimated in some cases at the very least 0.01 gsi wont be viable for all animals
	Some caveats with the results of the model.  First, growth is simulated as being extremely fast.  As previously stated, asymptotic size was reach by $\sim$10 days.  This is of course not representative of the real world, where individuals generally need several months to years to reach maturity.	
	The rapid growth may be due to several factors.  First is that metabolic cost may be underestimated.  Similar to \citeauthor{West2001} (\citeyear{West2001}), this study used resting metabolic rate to define metabolic costs.  However, this does not take other costs into account such as digestion and locomotion.
	%A rate which is typically used within energetic and growth studies, including early OGMs \parencite{Brett1971,Gillooly2001}  % ref for RMR in growth Tsoukali?  Burger? 
	This was addressed in traditional OGMs by \citeauthor{Hou2008} (\citeyear{Hou2008}).  However, due to the use of asymptotic mass in the parametrisation of this change, the same changes could not be used in this model.  Resting metabolic rate and active metabolic rates do not scale in the same way with mass \parencite{Gillooly2001, Weibel2004}.  The additional cost of active metabolic rate would cause a steeper scaling within the metabolic cost term, leading to more gradual growth (Fig. \ref{scaling_plot}). 
	As such, inclusion of active metabolic rates, while challenging to measure directly and implement, is needed.
	Additionally, there may be behavioural or physiological factors that would lead to an altered metabolic rate.  
	% TODO V presenting this as an issue but it is an area for improvement needs rephrasing and/or shifting
	In this regard, temperature plays a critical role.  It is well documented that a change in temperature will change many biological rates  (see \citeauthor{Peters1983} (\citeyear{Peters1983}), \citeauthor{Gillooly2001} (\citeyear{Gillooly2001}) etc.).  
	It has been shown that growth is dependent on temperature within fish. % sources
	For example, a 1.5\textdegree{}C increase in sea temperatures could result in a 15\% decrease in fish lengths \parencite{VanRijn2017}. % try to phrase better or bolster since it is lifted from diego
	The functional response data used in this study is standardised around 15\textdegree{}C \parencite{Pawar2012}.
	Meanwhile the metabolic cost is for an unspecified temperature \parencite{Peters1983}.  Using rates where the temperature effect is taken into account is crucial for model accuracy given the variation in rates that occurs over different temperatures.  Work such as \citeauthor{Barneche2014} (\citeyear{Barneche2014}) has investigated this effect, however the estimate for metabolic rate is several orders of magnitude lower than what was reported by \citeauthor{Peters1983} (\citeyear{Peters1983}), which is the rate used in this study.  Thus further investigation is required.  % should include a plot in the SI for this.
	
	In fish metabolic rate has also been shown to drop under starvation \parencite{Cook2000}.  In homeotherms feeding restriction has also been shown to also lower body temperature, since metabolic rate and core temperature are closely related in homeotherms \parencite{Ballor1991, Blanc2003,}.  
	
	Another possible point of error is the estimates for supply.  The parameters used from \citeauthor{Pawar2012} (\citeyear{Pawar2012}) are for a spectrum of animals from mammals to insects.  It is possible by reanalysing the data for only marine species, or more specifically only within taxon or species, predictions of supply could be improved \parencite{Marshall2019}.  % check dim paper for addressing rates in different environments
	
	A factor that is not taken into account in this model is that resources are not constant over time.  This can be implemented within the model by varying resource density over time.  The functional response will respond accordingly giving intake which varies through time.  One concern with implementing such a response is fluctuations are likely not experienced by all organisms in the same way.  For a fish with a small range a local fluctuation can be measured and described relatively simply.  However, for a fish with a very large range there is the possiblity leaving resource poor areas in search of richer waters.
	
	
	Despite these caveats, the patterns which arise from the model are promising.  The qualitative patterns seen in this study should not change even with slower growth.  This is because while the exact results are not representative of reality, the relationship between the values is.  For example, the manner in which supply and metabolic rate interact for different values of the metabolic scaling exponent.  The fact that a higher metabolic scaling exponent will cause intake and metabolic cost to intersect sooner does not change regardless of the absolute values.
	
 
	% However such field data is sparse. <- justify saying this some how. need a review or something.
	
	
	% What else can be done to address these issues
		% What is the minimum amount of GSI needed. reproductive structure will vary a bit between species and there is likely a minimum that a species must devote to these structures for them to be viable.
	
	
	% Emphasis that this is applicable to more than just fish - What needs to be done to apply the model to other animal groups.
	In conclusion, the model presented in this study is a promising base which can be expanded upon in a way that was not possible with previous OGMs allowing for much more controlled and detailed explanations of the factors controlling growth.  In contrast to all previous work, which assume optimal supply, the concept of varying supply is addressed using functional responses.  Additionally, qualitative evidence is provided supporting hyperallometric scaling in fish using energy budget as the basis.  The model can easily be applied to any animal taxon, not just fish, with some simple changes.  Additionally, there are clear directions to be explored to improve the model's accuracy.
	
	%TODO make sure that i mention that the increasing evidence for hyperallometric scaling in fish means that there is need to reconsider the stance of fishing mainly large fish and leaving behind the small ones.
	% TODO make sure that I spin the resources into conservation or management if possible?
	
	
	
	

%	\subsection{planned outline}
%	Overall summary of results and their meaning.
%	\\- I show that hyper allometry is possible within fish.  
%	\\- higher metabolic cost are leading to hyper allometry, even though I am seeing this for what the data considers too high (i.e. metabolic scaling expontent = 1) \parencite{Marshall2019} % see my google drive doc
%	\\- scaling in 3D will outscale metabolic cost in all cases Table \ref{parameters} because it is greater than 1, need reproductive cost to stop growth.
%	Shrinking and how it addresses the growth speed issue in this model to a degree in my opinion.
%	\\- Is like seeing the stage of growth we expect to see in a typical growth curve in reverse.  Where instead of growing into the asymptotic mass, it shrinks to it.  The energy balance remains the same.
%	\\
%	Issues with my model
%	- mass is too high, addressing a more detailed reproductive mass/metabolic rate may address this or reanalysing samraats data to be more taxon specific may do it.
%	\\
%	What else can be done.
%	\\- more fish specific intake using \cite{Pawar2012} methods as argued in \cite{Marshall2019}
%	\\- trying to model specific species with know growth curves to see how well our results map to them even if the curve isn't great it could be that reproductive output and final mass map well.
%	\\- more detailed metabolic cost since here is basically just RMR with higher exponents.  It may be that the constant shift too and this would lead to better predictions
%	\\- Incorporate oscillating resource density.  was attempted here but due to computational problems could not be implemented.  There is a lack of oscillation data this though and the goal would need to be far more specific since it will depend on scale and range of the organism.
%	\\- This same framework could be applied to other taxa with just a little adjustment (metabolic cost and reproductive equation?)
%	\\- There may be behavioural changes that occur at under feeding restriction that slows metabolism, similar to hibernation etc.
%	\\- temperature could be added here to adjust the metabolic and functional response rates %samraats data is standardised to 15C doesn't seem to be noted in peter's, hemingsen may show it
%	- develop the model to include storage like DEB or put this into a DEB framework.
%	\\
%	What is the minimum amount of GSI needed. reproductive structure will vary a bit between species and there is likely a minimum that a species must devote to these structures for them to be viable. %\cite{Gunderson1997} shows that mortality increases with higher GSI and larger inds have lower GSI does show low values for GSI (0.02) 
%	%GSI can reach 0.75 in invertebrates \cite{Parker2018}
%	\\
%	Emphasis that this is applicable to more than just fish - What needs to be done to apply the model to other animal groups. % need to open up gsi \cite{Parker2018} 0.75 inverts
%	\\
%	This is all energetically based, allowing us to ask questions of why or how there are exceptions to this in the case that organsims should be able to grow larger energetically speaking but may be limited genetically or via some other mechanism
	
	\nolinenumbers
%	\subsection{notes}
%	there are some integration issues that could allow for the appearance of more possible values
%	The shrinking helps to show what could be possible given there is a more gentle growth curve since the main difference is that the time when growth become more limited is just shown as shrinking vs slower growth.


	%%%%%%%%%% Conclusion %%%%%%%%%%
%\section{Conclusion}
%	\linenumbers
%	
%	
%	
%	\nolinenumbers
	
	%%%%%%%%%% Data Declaration %%%%%%%%%%
	\textbf{Code and Data Availability}

	Code is available at: \url{https://github.com/Don-Burns/Masters_Project}
	
	
	%%%%%%%%%% Bibliography %%%%%%%%%%
	\newpage
	\linenumbers
	
	%nicer indentation
	\addcontentsline{toc}{section}{\protect\numberline{}References}	
	
	\printbibliography
	
	% prob more official method of doing the above
%	\printbibliography[heading=bibintoc, title={References}]
	\nolinenumbers
	
	
	%%%%%%%%%% SI %%%%%%%%%%
	\newpage
	%to have table and figures numbered from 1 with prefix
\newcommand{\beginsupplement}{%
	\addcontentsline{toc}{section}{\protect\numberline{}Supplementary Information}
	\setcounter{table}{0}
	\renewcommand{\thetable}{S\arabic{table}}%
	\setcounter{figure}{0}
	\renewcommand{\thefigure}{S\arabic{figure}}%
}

\section*{Supplementary Information}
\beginsupplement

\subsection*{notes}
need section on value conversions and derivations

move any unreferenced sensitivity analyses here.


\begin{figure}
	\centering
	\includegraphics[width=\textwidth]{../../results/Sens_Maturation_Time_Long}
	\caption{Long time scale maturation time}
	\label{fig:sensmaturationtimelong}
\end{figure}

\begin{figure}
	\centering
	\includegraphics[width=\textwidth]{../../results/Sens_Maturation_Time_short}
	\caption{Short time scale maturation time}
	\label{fig:sensmaturationtimeshort}
\end{figure}

\begin{figure}
	\includegraphics[width=\textwidth]{../../results/Sens_Low_Resources_Maturation_Time_short}
	\caption{Low Resources Maturation}
	\label{fig:senslowresourcesmaturationtimeshort}
\end{figure}

\begin{figure}
	\centering
	\includegraphics[width=\textwidth]{../../results/Sens_Metabolic_Exponent}
	\caption{High Resources Metabolic Exponent}
	\label{fig:sensmetabolicexponent}
\end{figure}

\begin{figure}
	\centering
	\includegraphics[width=\textwidth]{../../results/Sens_Low_Resources_Metabolic_Exponent}
	\caption{Low Resource Metabolic Exponent}
	\label{fig:senslowresourcesmetabolicexponent}
\end{figure}

\begin{figure}
	\centering
	\includegraphics[width=\textwidth]{../../results/Sens_Resource_Density}
	\caption{Large scale Resource Resource Density}
	\label{fig:sensresourcedensity}
\end{figure}

\begin{figure}
	\centering
	\includegraphics[width=\textwidth]{../../results/Sens_Low_Resources_Resource_Density}
	\caption{Fine scale resource Density}
	\label{fig:senslowresourcesresourcedensity}
\end{figure}

\begin{figure}
	\centering
	\includegraphics[width=\textwidth]{../../results/Sens_Shrink}
	\caption{High Resources Shrinking}
	\label{fig:sensshrink}
\end{figure}

\begin{figure}
	\centering
	\includegraphics[width=\textwidth]{../../results/Sens_Low_Resources_Shrink}
	\caption{Low resources shrinking, missing values around 0 are because no viable solutions could be found.}
	\label{fig:senslowresourcesshrink}
\end{figure}




	


	
	
\end{document}